\documentclass[12pt,a4paper]{article}
\usepackage[left=3cm,right=2.5cm,top=2.5cm,bottom=2.5cm]{geometry}
\usepackage{fontspec}
\usepackage[main=thai,provide=*]{babel}
\babelprovide[main]{thai}
\babelfont{rm}[Path=docs/fonts/,Extension=.ttf,UprightFont=Sarabun-Regular,BoldFont=Sarabun-Bold,ItalicFont=Sarabun-Italic,BoldItalicFont=Sarabun-BoldItalic,Scale=1.06,Renderer=HarfBuzz]{Sarabun}
\babelfont{tt}[Path=docs/fonts/,Extension=.otf,UprightFont=CommitMono-400-Regular,BoldFont=CommitMono-700-Regular,ItalicFont=CommitMono-400-Italic,BoldItalicFont=CommitMono-700-Italic,Scale=MatchLowercase,Ligatures={TeXOff,NoContextual,NoDiscretionary},Renderer=HarfBuzz]{CommitMono}
\linespread{1.30}
\usepackage{booktabs,longtable,tabularx,array,graphicx,tikz,pgfplots,xcolor}
\usepackage{needspace}
% flafter forbids a float from rising above its declaration point on the page,
% which kept mid-paragraph text from being pushed below a figure declared later.
\usepackage{placeins}
\usepackage{flafter}
\usepackage{ragged2e}
\usepackage{framed}
\usepackage[titles]{tocloft}
% The build copy inserts this command only at dictionary-approved Thai word
% boundaries. Its flexible zero-width space supports full justification without
% allowing a Thai word to split between lines.
\DeclareRobustCommand{\thaiwordbreak}{\hspace{0pt plus 0.045em}}
\directlua{Babel.sea_enabled = false}
\hyphenpenalty=10000
\exhyphenpenalty=10000
\usetikzlibrary{positioning,arrows.meta,shapes.geometric}
\usepgfplotslibrary{statistics}
\pgfplotsset{compat=1.18}
\newcommand{\codebreaks}{\do\.\do\_\do\:\do\/\do\-\do\;}
\newcommand{\longcodebreaks}{\codebreaks\do\0\do\1\do\2\do\3\do\4\do\5\do\6\do\7\do\8\do\9\do\a\do\b\do\c\do\d\do\e\do\f\do\g\do\h\do\i\do\j\do\k\do\l\do\m\do\n\do\o\do\p\do\q\do\r\do\s\do\t\do\u\do\v\do\w\do\x\do\y\do\z\do\A\do\B\do\C\do\D\do\E\do\F\do\G\do\H\do\I\do\J\do\K\do\L\do\M\do\N\do\O\do\P\do\Q\do\R\do\S\do\T\do\U\do\V\do\W\do\X\do\Y\do\Z}
\newcommand{\code}[1]{\bgroup\def\UrlBreaks{\codebreaks}\path{#1}\egroup}
\newcommand{\codeL}[1]{\bgroup\def\UrlBreaks{\longcodebreaks}\path{#1}\egroup}
% \codeS breaks a semicolon-joined path list only between entries, never inside one.
\newcommand{\codeS}[1]{\bgroup\def\UrlBreaks{\do\;}\path{#1}\egroup}
% \cmd keeps visible spaces in shell invocations: unlike \code it never lets
% the space between the command and its argument collapse during typesetting.
\newcommand{\cmd}[1]{\texttt{#1}}
\usepackage{url}
\usepackage[unicode,colorlinks=true,linkcolor=black,urlcolor=blue,citecolor=black]{hyperref}
\usepackage{caption}
\captionsetup{font=small,labelfont=bf,justification=raggedright,singlelinecheck=false}
\urlstyle{tt}
\clubpenalty=10000
\widowpenalty=10000
\displaywidowpenalty=10000
\pdfstringdefDisableCommands{\def\thaiwordbreak{}}
\hypersetup{pdftitle={รายงานทางเทคนิค ระบบจัดตารางสอบ},pdfauthor={โครงการระบบจัดตารางสอบ},pdflang={th},pdfsubject={รายงานทางเทคนิค: การนำเข้าข้อมูล การค้นหาตารางสอบ และการตรวจสอบผลลัพธ์},pdfkeywords={จัดตารางสอบ, ระบบจัดตารางสอบ, รายงานทางเทคนิค, ไทย}}
\renewcommand{\cftdotsep}{2.5}
\renewcommand{\cftsecleader}{\cftdotfill{\cftdotsep}}
\renewcommand{\cftsubsecleader}{\cftdotfill{\cftdotsep}}
\renewcommand{\cftsubsubsecleader}{\cftdotfill{\cftdotsep}}
\renewcommand{\cftfigleader}{\cftdotfill{\cftdotsep}}
\renewcommand{\cfttableader}{\cftdotfill{\cftdotsep}}
\setlength{\cftbeforesecskip}{5pt}
\hypersetup{bookmarksdepth=3}
\setlength{\parindent}{1.1cm}
\setlength{\parskip}{5pt}
\setlength{\emergencystretch}{3em}
\setlength{\tabcolsep}{4pt}
% All table forms share the same centered placement.  Longtable uses its own
% left/right glue, so set it explicitly instead of relying on class defaults.
\setlength{\LTleft}{\fill}
\setlength{\LTright}{\fill}
\renewcommand{\arraystretch}{1.45}
\renewcommand{\topfraction}{0.9}
\renewcommand{\bottomfraction}{0.8}
\renewcommand{\textfraction}{0.08}
\renewcommand{\floatpagefraction}{0.75}
\setcounter{topnumber}{3}
\setcounter{bottomnumber}{2}
\setcounter{totalnumber}{5}
\newcolumntype{L}[1]{>{\raggedright\arraybackslash}p{#1}}
\newcolumntype{R}[1]{>{\raggedleft\arraybackslash}p{#1}}
\newcolumntype{Y}{>{\RaggedRight\arraybackslash}X}
\newcommand{\cmdcell}[1]{\mbox{\cmd{#1}}}
\justifying
\setlength{\parindent}{1.1cm}
% Keep explicit spaces around English terms from absorbing most of a justified
% line; the remaining adjustment is shared by Thai dictionary boundaries.
\spaceskip=0.25em plus 0.05em minus 0.04em
\xspaceskip=0.25em plus 0.05em minus 0.04em
\renewcommand{\texttt}[1]{\textbf{#1}}
\newcommand{\file}[1]{\texttt{#1}}
\newcommand{\course}[1]{\texttt{#1}}

\newcounter{chapter}
\newcommand{\chapter}[1]{%
  \clearpage
  % Keep the chapter's opening page free of top floats, so no figure may be
  % placed above the chapter number and title on their page.
  \suppressfloats[t]%
  \needspace{12\baselineskip}%
  \refstepcounter{chapter}%
  \refstepcounter{section}%
  \phantomsection
  \addcontentsline{toc}{section}{\protect\numberline{\thesection}#1}%
  \markboth{#1}{#1}%
  \noindent{\Large\bfseries \thesection\par}%
  \vspace{0.3em}%
  \noindent{\Huge\bfseries\RaggedRight #1\par}%
  \vspace{1.1em}%
}
\counterwithin{figure}{chapter}
\counterwithin{table}{chapter}
\tikzset{box/.style={draw,rounded corners,align=center,text width=4.5cm,minimum height=1cm},arr/.style={-{Stealth},thick}}
\tikzset{every node/.append style={execute at begin node={\linespread{1}\selectfont}}}
\pgfplotsset{colormap={demand}{rgb255(0)=(246,248,250) rgb255(1)=(0,109,44)}}
% One uniform box style: size labels are on the x axis, so color carries no meaning.
\pgfplotsset{boxplot/every box/.style={draw=black,fill=black!12},
  boxplot/every whisker/.style={draw=black},
  boxplot/every median/.style={draw=black},
  boxplot/every outlier/.style={draw=black,fill=black}}
\begin{document}
\begin{titlepage}
\centering
\vspace*{1cm}
\begin{tabular}{@{}c@{\hspace{1.2cm}}c@{}}

\includegraphics[height=2.2cm,keepaspectratio]{docs/usage-report/branding/fitm-logo.png} &

\includegraphics[height=2.2cm,keepaspectratio]{docs/usage-report/branding/kmutnb-logo.png}
\end{tabular}\par
\vspace{0.5cm}
{\large คณะเทคโนโลยีอุตสาหกรรมและการจัดการ\par
มหาวิทยาลัยเทคโนโลยีพระจอมเกล้าพระนครเหนือ\par}
\vspace{0.9cm}
{\Huge รายงานทางเทคนิค\par}
\vspace{0.7cm}{\Huge\bfseries ระบบจัดตารางสอบ Exam Scheduler\par}
\vspace{1cm}{\Large การนำเข้าข้อมูล การค้นหาตารางสอบ\par และการตรวจสอบผลลัพธ์\par}
\vspace{1.2cm}{\large 6 กันยายน พ.ศ. 2569\par}
\end{titlepage}
\pagenumbering{gobble}
\section*{บทคัดย่อ}
ระบบจัดตารางสอบ เป็นระบบที่รับข้อมูลรายวิชา กฎการจัดสอบ และข้อมูลห้องจากแฟ้มต้นทาง แล้วแปลงเป็นรายการสอบเพื่อค้นหาวันและเวลาที่สอดคล้องกับข้อกำหนด รายงานนี้อธิบายแบบจำลองและลำดับการประมวลผลด้วยชื่อฟังก์ชันที่ตรวจสอบได้ โดยแยกข้อบังคับออกจากเกณฑ์ที่ใช้ปรับปรุงคุณภาพตาราง

รายงานทดสอบระบบด้วยชุดข้อมูลสังเคราะห์ที่สร้างซ้ำได้สามขนาด คือ 41, 100 และ 189 วิชา ครอบคลุม 8 สถานการณ์ที่กำหนดผลคาดหมายไว้ล่วงหน้า ทั้งกรณีใช้งานปกติและกรณีผิดพลาด การทดสอบเรียกโปรแกรมผ่านบรรทัดคำสั่ง ตรวจผลซ้ำด้วยเครื่องมือตรวจยืนยันที่แยกจากกระบวนการค้นหา และบันทึกภาพหน้าจอส่วนติดต่อผู้ใช้เป็นหลักฐานประกอบ

ในชุดข้อมูลหลัก ระบบกำหนดวันเวลาให้รายการสอบที่รับผิดชอบครบ 82 จาก 82 รายการ และสถานการณ์ทดสอบทั้งแปดให้ผลตรงกับที่กำหนดไว้ล่วงหน้า การค้นหาพิสูจน์ได้ว่าจำนวนรายการที่จัดได้ถึงขอบเขตบนที่คำนวณได้ในทั้งสามขนาด ส่วนจำนวนรายการในวันที่ควรหลีกเลี่ยงพิสูจน์ว่าต่ำที่สุดได้ในชุดเล็กและชุดกลาง เวลาทำงานครบกระบวนการผ่าน CLI มีค่ามัธยฐานประมาณ 0.56, 0.99 และ 2.12 วินาทีตามลำดับขนาด และการทดสอบความอ่อนไหวของพารามิเตอร์ของฟังก์ชันต้นทุนไม่พบความไวสูงของความครอบคลุมภายในช่วงค่าที่ทดลอง

ผลทั้งหมดได้จากข้อมูลสังเคราะห์ที่เลือกเงื่อนไขให้จัดครบได้ จึงไม่ใช่ตัวอย่างแบบไม่มีอคติของความยากทั้งหมด ผลด้านเวลาใช้ได้เฉพาะภายใต้สภาพแวดล้อมที่บันทึกไว้ และขอบเขตของข้อสรุปทั้งหมดสรุปไว้เป็นตารางการยืนยันในบทที่ 1
\clearpage\pagenumbering{roman}
\begingroup\setlength{\parskip}{0pt}\tableofcontents\endgroup
\clearpage
\listoffigures
\clearpage
\listoftables
\clearpage\pagenumbering{arabic}\setcounter{page}{1}
\chapter{บทนำและขอบเขต}
\subsection{ปัญหาและวัตถุประสงค์}
การจัดตารางสอบต้องประสานข้อมูลที่มีหน่วยและความหมายต่างกัน รายวิชาหนึ่งอาจมีหลายตอนเรียน นักศึกษากลุ่มเดียวกันอาจเรียนหลายวิชา และบางรายวิชามีวันเวลาสอบระบุมาแล้ว การย้ายรายการใดรายการหนึ่งจึงอาจกระทบรายการอื่น แม้ตารางที่มองเห็นจะมีช่องว่างอยู่ก็ตาม ปัญหาสำคัญไม่ใช่เพียงการหาช่องว่าง แต่เป็นการยืนยันว่าช่องนั้นใช้ได้ภายใต้ข้อมูลและกฎทั้งหมดที่เกี่ยวข้อง

รายงานนี้มีวัตถุประสงค์สามประการ ประการแรก อธิบายแบบจำลองและลำดับการประมวลผลด้วยชื่อฟังก์ชันที่ตรวจสอบได้ ประการที่สอง แสดงตัวอย่างที่ทำงานได้จริงทั้งกรณีปกติและกรณีผิดพลาด ประการที่สาม กำหนดขอบเขตของผลการทดสอบให้ชัดเจน โดยไม่ใช้การค้นหาไม่สำเร็จเป็นหลักฐานว่าปัญหาไม่มีคำตอบ
\subsection{หน่วยของข้อมูล}
คำว่า รายวิชา หมายถึงระเบียนที่ระบุด้วยรหัสวิชา ส่วน ตอนเรียน เป็นข้อมูลย่อยที่เก็บกลุ่มนักศึกษาและรายละเอียดจากต้นทาง รายการสอบ ผูกกับรายวิชาและประเภทการสอบ ได้แก่กลางภาคหรือปลายภาค จำนวนตอนเรียนจึงไม่ใช่จำนวนรายการสอบโดยตรง รายวิชา 41 วิชาในชุดหลักให้รายการสอบที่ระบบรับผิดชอบ 82 รายการ

รายวิชาที่มีบทบาท \code{managed} อยู่ในขอบเขตที่ระบบจัดตาราง ส่วน \code{context} ให้ข้อมูลประกอบ โดยเฉพาะรายการที่มีวันเวลาแล้วและอาจชนกับกลุ่มนักศึกษาเดียวกัน การไม่นับรายวิชาประกอบเป็นผลงานที่ระบบจัดใหม่ช่วยให้ตัวเลขความครอบคลุมไม่สูงเกินจริง
\subsection{อัตลักษณ์ของซอฟต์แวร์และการทำซ้ำ}\label{ch:identity}
ซอฟต์แวร์ที่ศึกษาคือ ระบบจัดตารางสอบ แพ็กเกจ \code{exam-scheduler-core} รุ่น 0.1.0 เผยแพร่ที่ \code{https://github.com/ItsBreadTime/examscheduler} ภายใต้สัญญาอนุญาต MIT การอ้างอิงสถานะโปรแกรมใช้ Git revision ของ repository ร่วมกับ snapshot ที่บันทึกโดยเครื่องมือ \code{docs/repro/snapshot.ts} ซึ่งเก็บแฮชของทุกแฟ้มที่อยู่นอก \code{docs/} การเปรียบเทียบก่อนและหลังการทำรายงานจึงตรวจได้ว่าโค้ดที่หลักฐานอ้างถึงไม่ถูกเปลี่ยนระหว่างการทำรายงาน โดยมีข้อยกเว้นที่บันทึกไว้เฉพาะบันทึกการทำงานชั่วคราวของเครื่องมืออัตโนมัติ หัวข้อในรายงานที่อ้างอิงแฟ้ม implementation รวบรวมไว้ในภาคผนวก \ref{app:trace}

รายงานนี้ใช้สถานะซอฟต์แวร์สองจุดเป็นหลักฐานอ้างอิง หลักฐานการทำงานผ่านบรรทัดคำสั่งและการวัดเวลาสร้างใหม่ที่ revision ปัจจุบันเมื่อวันที่ 6 กันยายน 2569 ส่วนภาพหน้าจอของส่วนติดต่อผู้ใช้บันทึกไว้ตั้งแต่จุดอ้างอิงเดิมของวันที่ 5 กันยายน 2569 ซึ่งต่อมาถูกบันทึกเข้าประวัติของ repository แล้ว การเปรียบเทียบแฮชรายแฟ้มระหว่างสองจุดอ้างอิงยืนยันว่าโค้ดของกระบวนการจัดตารางทั้งหมดภายใต้ \code{src/lib} ไม่ต่างกันแม้แต่ไบต์เดียว สิ่งที่เปลี่ยนคือแฟ้มระดับส่วนติดต่อผู้ใช้และการเผยแพร่ ซึ่งไม่มีผลต่อผลลัพธ์การจัดตารางที่รายงานอ้างอิง

รายงานฉบับนี้เพิ่มจุดอ้างอิงที่สามสำหรับชั้น cloud คือการตรวจระบบที่เผยแพร่จริงที่ \code{https://schedule.breadtm.xyz/} เมื่อวันที่ 6 กันยายน 2569 แฟ้ม JavaScript CSS และนโยบายความปลอดภัยของเนื้อหาที่เซิร์ฟเวอร์จริงตอบกลับตรงกับผลการสร้างใหม่จาก revision ปัจจุบันทุกแฟ้ม ผลการตรวจอยู่ในแฟ้ม \code{docs/data/cloud-verification.json} และรายละเอียดการบันทึกภาพหน้าจออยู่ในแฟ้ม \code{docs/screenshots/cloud-screenshots.json} ข้ออ้างเกี่ยวกับระบบที่เผยแพร่จึงอ้างถึงโค้ดชุดเดียวกับหลักฐานอื่นของรายงาน

สถานะซอฟต์แวร์ที่รายงานอ้างอิงมีดังนี้
\begingroup\def\UrlBreaks{\longcodebreaks}\input{docs/.build/reproducibility.tex}\endgroup
\subsection{ศัพท์ที่ใช้ตลอดรายงาน}
\begin{longtable}{L{5.2cm}L{8cm}}
\caption{ศัพท์ที่ใช้ตลอดรายงาน}\label{tab:terms}\\
\toprule คำในรายงาน & ความหมาย \\\midrule
หน่วยสอบพร้อมกัน (same-time unit) & กลุ่มรายการสอบที่ต้องเลือกวันเวลาเดียวกัน \\
เลือกตัวแปรที่มีทางเลือกน้อยที่สุด (MRV) & เลือกหน่วยจากจำนวนช่องเวลาที่ใช้ได้ ณ สถานะค้นหาปัจจุบัน \\
ต้นทุนแบบเรียงลำดับความสำคัญ (lexicographic cost) & เปรียบเทียบองค์ประกอบจากซ้ายไปขวา โดยองค์ประกอบแรกที่ต่างกันเป็นตัวตัดสิน \\
กฎอ้างอิงล้าสมัย (stale rule) & กฎอ้างถึงวิชาที่ไม่มีในขอบเขตข้อมูลที่ใช้ \\
กราฟข้อขัดแย้ง (conflict graph) & โครงสร้างเชื่อมหน่วยที่ไม่สามารถสอบทับซ้อนกัน \\
งบการค้นหา (search budget) & ขอบเขตงานค้นหาที่ตั้งไว้ ไม่ใช่ระยะเวลาเป็นวินาที \\
การตัดกิ่งด้วยขอบเขตบน (upper-bound pruning) & งดสำรวจสาขาที่ไม่อาจเพิ่มความครอบคลุมเหนือผลที่มี \\\bottomrule
\end{longtable}
คำศัพท์เหล่านี้ใช้แยกระหว่างกลไกกับข้อสรุป เช่น งบการค้นหาหมดเป็นสถานะของกระบวนการ แต่ไม่ใช่ข้อสรุปว่ารายวิชาทั้งหมดจัดสอบไม่ได้
\subsection{สิ่งที่รายงานนี้ยืนยันและไม่ยืนยัน}\label{ch:claims}
ตารางนี้สรุปข้ออ้างของรายงานทั้งฉบับ เพื่อให้การอ่านซ้ำในบทสรุปอ้างจุดเดียวกัน รายงานตั้งคำถามประเมินผลห้าข้อ ได้แก่ E1 ระบบจัดรายการ managed ได้ครบหรือไม่ E2 ข้อกำหนดหลัก ได้แก่ การไม่ทับซ้อนของกลุ่มนักศึกษา กฎสอบพร้อมกัน และวันเวลาต้นทาง ถูกตรวจซ้ำหรือไม่ E3 จำนวนรายการที่จัดได้ถึงขอบเขตบนที่พิสูจน์ได้หรือไม่ E4 จำนวนรายการในวันที่ควรหลีกเลี่ยงต่ำที่สุดหรือไม่ และ E5 ส่วนติดต่อผู้ใช้กับ cloud อยู่ในขอบเขตการทดสอบเพียงใด แต่ละแถวระบุคำถามที่ตอบ หลักฐานที่ใช้ และสิ่งที่ข้ออ้างนั้นไม่ครอบคลุม บทต่อ ๆ ไปขยายรายละเอียดของแต่ละแถวโดยไม่ต้องซ้ำข้อสรุปเดิมทั้งย่อหน้า

\begin{center}\begin{minipage}{\textwidth}\centering
\captionof{table}[ขอบเขตข้ออ้างหลักของรายงาน]{ขอบเขตของข้ออ้างหลักของรายงาน: แต่ละแถวระบุคำถามประเมินผลที่ตอบ หลักฐานที่ใช้ และขอบเขตที่ข้ออ้างไม่ครอบคลุม}\label{tab:claims}
\begin{tabularx}{\textwidth}{@{}L{0.30\textwidth}L{0.26\textwidth}Y@{}}\toprule
ข้ออ้างที่ยืนยัน & หลักฐานและตำแหน่ง & ขอบเขตที่ไม่ยืนยัน \\\midrule
E1 รายการสอบ managed ของชุดหลักมีวันเวลาครบ 82 รายการ (ค้นหาใหม่ 78 รายการ) & ผล CLI บทที่ \ref{ch:function} & ไม่ยืนยันว่าคุณภาพทุกด้านเหมาะที่สุด \\
E2 การจัดวันเวลาไม่ทับซ้อนกับกลุ่มนักศึกษาที่กำหนด คู่วิชาตามกฎสอบพร้อมกันอยู่ช่วงเดียวกัน และวันเวลาต้นทางถูกรักษาไว้ & ผลการตรวจซ้ำ บทที่ \ref{ch:method} & ครอบคลุมเฉพาะกลุ่มที่มีในข้อมูล \\
E3 จำนวนที่จัดได้ถึงขอบเขตบนที่พิสูจน์ได้ & สถิติการค้นหา บทที่ \ref{ch:perf} & ไม่ยืนยันความเหมาะสมที่สุดขององค์ประกอบคุณภาพอื่น \\
E4 จำนวนรายการในวันที่ควรหลีกเลี่ยงต่ำที่สุด & สถิติการค้นหา ยืนยันในชุดเล็กและกลาง ไม่ยืนยันในชุดใหญ่ & รายละเอียดอยู่ในบทที่ \ref{ch:perf} \\
E5 ความสามารถของส่วนติดต่อผู้ใช้มุมมองหลักเจ็ดมุมมอง และการบันทึกกับแบ่งปันโครงการบนระบบที่เผยแพร่ & ภาพหน้าจอจากจุดอ้างอิง 5 กันยายน และการตรวจระบบที่เผยแพร่ 6 กันยายน บทที่ \ref{ch:method} & ไม่ครอบคลุมการโต้ตอบที่ลึกกว่าการเปิดดู เช่น การแก้ตารางด้วยมือ \\\bottomrule
\end{tabularx}\end{minipage}\end{center}
\chapter{สถาปัตยกรรมและแบบจำลองข้อมูล}\label{ch:arch}
\subsection{องค์ประกอบของระบบ}
ระบบมีช่องทางเรียกใช้สองแบบที่ใช้ไลบรารีแกนร่วมกัน ช่องทางแรกคือส่วนติดต่อผู้ใช้บนเบราว์เซอร์ ซึ่งส่งงานหนักไปยัง Web Worker ผ่าน \code{runJob} โดยสร้าง worker ต่องานและยุติเมื่อสำเร็จ ผิดพลาด หรือถูกยกเลิก ช่องทางที่สองคือ CLI สองคำสั่ง คือ \code{schedule} ที่ทำงานครบตั้งแต่นำเข้าจนเขียนผลลัพธ์ และ \code{validate} ที่ตรวจโครงการที่ส่งออกแล้วซ้ำ การประมวลผลจริงทั้งการค้นหาวันเวลาและการตรวจซ้ำอยู่ในไลบรารีแกน จึงเป็นหลักฐานร่วมของทุกช่องทางเรียกใช้ ชั้น cloud รองรับการเก็บ revision และสิทธิ์การเข้าถึง โดยการบันทึกเมื่อ revision ไม่ตรงกันตอบกลับด้วย \code{REVISION_CONFLICT} ระบบที่เผยแพร่จริงผ่านการตรวจตามหัวข้อ \ref{sec:cloud}

\begin{figure}[tbp]\centering
\begin{tikzpicture}[node distance=0.5cm]
\node[box,text width=4.7cm] (ui) {ส่วนติดต่อผู้ใช้ (Svelte)};
\node[box,right=1.6cm of ui,text width=4.7cm] (worker) {Web Worker\\ \code{runJob} ต่องานหนึ่งครั้ง};
\node[box,below=of ui,text width=4.7cm] (cli) {CLI \code{schedule} และ \code{validate}};
\node[box,below=1.1cm of worker,text width=4.7cm] (core) {ไลบรารีแกน \code{src/lib}\\ นำเข้า ค้นหา ตรวจซ้ำ};
\node[box,below=1.7cm of cli,text width=4.7cm] (out) {ผลลัพธ์ JSON หลักฐาน audit};
\node[box,right=1.6cm of out,text width=4.7cm] (cloud) {ชั้น cloud (worker API)\\ ตรวจบนระบบที่เผยแพร่แล้ว};
\draw[arr](ui) to (worker);
\draw[arr](cli) to (core);
\draw[arr](worker) to (core);
\draw[arr](core) to (out);
\draw[arr](ui.east) -- ++(0.8cm,0) |- (cloud.north);
\end{tikzpicture}
\caption[องค์ประกอบของระบบ]{องค์ประกอบของระบบ เส้นลูกศรแสดงทิศทางการทำงาน โดยหลักฐานทดสอบ CLI และการแสดงผล UI แยกกันตามตารางที่ \ref{tab:capability} ช่องทาง cloud ตรวจบนระบบที่เผยแพร่จริงตามหัวข้อ \ref{sec:cloud}}
\label{fig:arch}
\end{figure}

สถานะความสามารถของระบบตามขอบเขตการทดสอบสรุปได้ดังนี้

\begin{center}\begin{minipage}{\textwidth}\centering
\needspace{14\baselineskip}\captionof{table}{สถานะความสามารถของระบบตามขอบเขตของรายงาน}\label{tab:capability}
\begin{tabularx}{\textwidth}{@{}L{0.24\textwidth}L{0.30\textwidth}Y@{}}\toprule
ความสามารถ & สถานะ & หลักฐานและขอบเขต \\\midrule
นำเข้าและจัดตารางผ่าน CLI & ทดสอบแล้ว & ชุดปกติและชุดข้อผิดพลาด \\
ตรวจต้นทางด้วย \code{inspect} & ทดสอบแล้ว & กรณีปฏิเสธข้อมูลสี่ชุด \\
\cmd{npm run validate} & มี implementation แต่ยังไม่ทดสอบในรายงาน & รับ Project JSON ที่ส่งออก ไม่รับ manifest \\
การแสดงผลใน UI & ทดสอบแล้ว & ภาพหน้าจอมุมมองหลักในบทที่ \ref{ch:method} \\
เส้นทาง Web Worker & หลักฐานทางอ้อม & ซอร์สโค้ดและ workflow ของ UI ยังไม่มีบันทึกการทำงานแยกเฉพาะ worker \\
การเก็บและแบ่งปันโครงการบน cloud & ทดสอบแล้ว & สัญญา API และการไหลบันทึกและแชร์บนระบบที่เผยแพร่จริง หัวข้อ \ref{sec:cloud} \\
ปฏิทินวันหยุดอัตโนมัติ & ยังไม่มี & ปัจจุบันต้องกำหนดวันที่ใน settings \\\bottomrule
\end{tabularx}\end{minipage}\end{center}
ความสามารถอื่นที่ไม่ปรากฏในตารางนี้อยู่นอกขอบเขตของรายงานฉบับนี้ สถานะ ทดสอบแล้ว หมายถึงมี implementation และถูกทดสอบในรายงานนี้ มี implementation แต่ยังไม่ทดสอบในรายงาน หมายถึงมีโค้ดแต่หลักฐานในรายงานไม่ครอบคลุม และ ยังไม่มี หมายถึงไม่มี implementation

\cmd{npm run validate} เรียก \code{validateProject} และสร้าง audit จากสถานะโครงการที่มีอยู่ จึงเหมาะกับการตรวจซ้ำหลังส่งออก ส่วน \code{schedule} ในโหมด \code{inspect} ตรวจโครงการที่เพิ่งนำเข้าจากแฟ้มต้นทาง ทั้งสองใช้ตัวตรวจเดียวกันบางส่วนแต่มีจุดเริ่มต้นต่างกัน
\FloatBarrier
\subsection{จากแฟ้มต้นทางสู่โครงการ}
การประมวลผลเริ่มจาก manifest ซึ่งบอกตำแหน่งแฟ้ม บทบาทข้อมูล และการตั้งเวลา เส้นทางแฟ้มใน manifest อ้างอิงจากโฟลเดอร์ที่เก็บ manifest จึงย้ายชุดข้อมูลทั้งโฟลเดอร์ได้โดยไม่ต้องแก้ตำแหน่งแบบสัมบูรณ์ ขั้นนำเข้ารวมระเบียนเป็นโครงการและตรวจปัญหาก่อนค้นหา ผลการตรวจขั้นนี้ต้องอ่านแยกจากผลตรวจหลังจัดตาราง

การค้นหาใช้ข้อมูลที่ผ่านการแปลงแล้ว ไม่อ่านตารางคำนวณซ้ำทุกครั้งที่ลองช่อง

\code{readSpreadsheet} อ่านแฟ้มและคำนวณ SHA-256 จากไบต์ต้นฉบับ ข้อมูลที่มาประกอบด้วยชื่อแฟ้ม ชนิดข้อมูล ขนาด และเวลานำเข้า การอ้างอิงแหล่งข้อมูลในระเบียนระบุแผ่นงาน แถว คอลัมน์ และชื่อเขตข้อมูลได้ จึงตามกลับไปยังตำแหน่งที่ต้องแก้ในต้นทาง อย่างไรก็ตาม แฮชยืนยันความเหมือนของเนื้อหา ไม่ได้ยืนยันผู้จัดทำหรือความถูกต้องของต้นฉบับ

\code{importProject} รวมข้อมูล ตรวจแหล่งซ้ำและบทบาทที่ขัดกัน แล้วสร้างรายการสอบผ่าน \code{buildEvents} หากวิชาเดียวกันปรากฏทั้งในข้อมูล managed และ context ระบบรายงาน \code{SOURCE_ROLE_CONFLICT} แทนที่จะเลือกบทบาทใดบทบาทหนึ่งโดยไม่แจ้งผู้ใช้ การแยกบทบาทมีผลทั้งต่อจำนวนรายการที่ต้องจัดและต่อการคิดภาระการสอบ

\begin{figure}[tbp]\centering
\begin{tikzpicture}[node distance=0.45cm]
\node[box] (a) {แฟ้มรายวิชา กฎ และห้อง};
\node[box,below=of a] (b) {นำเข้าและตรวจข้อมูลต้นทาง};
\node[box,below=of b] (c) {สร้างรายการสอบและหน่วยค้นหา};
\node[box,below=of c] (d) {ค้นหาวันเวลา};
\node[box,below=of d] (f) {ตรวจซ้ำและสร้างหลักฐานผลลัพธ์};
\foreach \x/\y in {a/b,b/c,c/d,d/f}\draw[arr](\x) to (\y);
\end{tikzpicture}
\caption{ข้อมูลผ่านขั้นใดบ้างก่อนเป็นผลลัพธ์ของ CLI}
\label{fig:pipeline}
\end{figure}
\FloatBarrier
\subsection{จากตอนเรียนสู่รายการสอบ}
วันเวลาจากต้นทางมีลำดับความสำคัญเหนือค่าล็อกที่กำหนดภายหลัง รายการที่มีวันเวลาต้นทางไม่ตรงกันระหว่างตอนเรียน หรือมีข้อมูลเวลาผิดรูปแบบ ระบบจะไม่สร้างเวลาใหม่เพื่อกลบความผิดปกติของข้อมูลต้นทาง

กฎสอบพร้อมกันเชื่อมโยงเป็นองค์ประกอบต่อเนื่องก่อนเลือกเฉพาะรายวิชาที่ยังใช้งาน ทำให้รหัสกลางที่ล้าสมัยไม่ตัดความสัมพันธ์ต่อเนื่องของกฎ

หน้าจอข้อมูลและการตั้งค่าเป็นจุดเริ่มต้นของการทำงานในเบราว์เซอร์ ผู้ใช้เพิ่มแฟ้มรายวิชา เลือกบทบาทของแต่ละแฟ้ม ระบุเงื่อนไขการสอบ และตั้งช่วงสอบกลางภาคกับปลายภาค ภาพถัดไปสรุปสิ่งที่หน้าจอตรวจสอบข้อมูลแสดงหลังการนำเข้า และแผนภาพท้ายหัวข้ออธิบายการแปลงแถวต้นทางเป็นรายการสอบ
\begin{figure}[tbp]\centering
\includegraphics[width=\textwidth]{docs/screenshots/sources.png}
\caption[หน้าจอข้อมูลและการตั้งค่า]{หน้าจอข้อมูลและการตั้งค่าจากชุดข้อมูลหลัก แสดงการเพิ่มแฟ้มและกำหนดช่วงสอบ}
\label{fig:sources}
\end{figure}
\begin{figure}[tbp]\centering
\includegraphics[width=\textwidth]{docs/screenshots/review.png}
\caption[หน้าจอตรวจสอบข้อมูล]{หน้าจอตรวจสอบข้อมูลจากชุดข้อมูลหลัก แสดงแฟ้มที่นำเข้า แฮชที่มา และสรุปช่วงสอบ}
\label{fig:review}
\end{figure}
\begin{figure}[tbp]\centering
\begin{tikzpicture}[node distance=0.6cm]
\node[box] (a) {แถวต้นทางหลายตอนเรียน};
\node[box,below=of a] (b) {รายวิชาและสหภาพกลุ่มนักศึกษา};
\node[box,below left=0.7cm and 0.1cm of b,text width=4cm] (c) {รายการสอบกลางภาค};
\node[box,below right=0.7cm and 0.1cm of b,text width=4cm] (d) {รายการสอบปลายภาค};
\draw[arr](a) to (b);\draw[arr](b) to (c);\draw[arr](b) to (d);
\end{tikzpicture}
\caption{แถวหลายตอนเรียนแปลงเป็นรายการสอบอย่างไร}
\label{fig:transform}
\end{figure}
\subsection{การแทนวันเวลา}\label{ch:time}
\code{ExamTiming} เก็บวันปฏิทินแบบ ISO ร่วมกับนาทีเริ่มและสิ้นสุดภายในวัน ไม่เก็บ instant หรือ timezone offset ค่านาทีเป็นจำนวนเต็ม โดยเริ่มตั้งแต่ 0 สิ้นสุดไม่เกิน 1440 และเริ่มน้อยกว่าสิ้นสุด การทดลองกำหนด timezone เป็น Asia/Bangkok เพื่อให้สภาพแวดล้อมสม่ำเสมอ ส่วนการคำนวณวันในสัปดาห์ใช้ UTC จากวันปฏิทินเพื่อหลีกเลี่ยงการเลื่อนวันตามเขตเวลาของเครื่อง

ฟังก์ชัน \code{overlaps} ตรวจว่าจุดเริ่มที่มากที่สุดน้อยกว่าจุดสิ้นสุดที่น้อยที่สุดอย่างเคร่งครัด และ \code{timingsOverlap} กำหนดให้วันตรงกันด้วย ช่วง 09:00 ถึง 12:00 กับ 12:00 ถึง 15:00 จึงไม่ทับซ้อน ระบบไม่เพิ่มเวลาพักหรือเวลาเตรียมห้องโดยอัตโนมัติ การอ่านวันที่จากแฟ้มต้นทางอธิบายในภาคผนวก \ref{app:dates}

\chapter{อัลกอริทึมการจัดตาราง}\label{ch:solver}
\subsection{สัญญาของปัญหาการค้นหา}
ก่อนอธิบายรายละเอียด ตารางนี้กำหนดปัญหาที่ \code{solveUnits} แก้ในรูปแบบมาตรฐานของการค้นหาแบบมีข้อกำหนด ทุกแถวตรวจสอบได้จากซอร์สโค้ดตามภาคผนวก \ref{app:trace}

\begin{longtable}{L{3.2cm}L{9.8cm}}\toprule
องค์ประกอบ & คำนิยาม \\\midrule
ตัวแปร & หน่วยสอบพร้อมกัน ได้แก่รายการเดี่ยว หรือกลุ่มรายการที่กฎผูกให้สอบพร้อมกัน ตัวแปรหนึ่งตัวได้เวลาเดียว \\
โดเมนของค่า & ผลคูณของวันในช่วงสอบกับช่วงเวลารายวันที่ใช้ได้ กรองด้วยข้อบังคับด้านล่าง \\
ข้อบังคับแบบแข็ง & วันเวลาต้นทางของตอนเรียนที่มีอยู่ กลุ่มนักศึกษาไม่ทับซ้อน ขอบเขตวันของรายการ นโยบายปฏิทิน \code{never} ความหมายเต็มวัน และกฎสอบพร้อมกัน \\
เป้าหมายอันดับหนึ่ง & เพิ่มน้ำหนักรวมของรายการที่จัดได้ (จำนวนรายการที่ยังไม่มีเวลาในหน่วย) \\
เป้าหมายรองตามลำดับ & ลดจำนวนรายการในวันที่ควรหลีกเลี่ยง ลดภาระกลุ่มนักศึกษา ลดความต้องการใช้ห้องพร้อมกันในช่วงเวลาเดียวกัน \\
การหยุดค้นหา & พิสูจน์ครอบคลุมถึงขอบเขตบน หรือใช้งบการค้นหาครบตามที่ตั้ง \\\bottomrule
\end{longtable}
\subsection{หน่วยค้นหา ช่องเวลา และกราฟข้อขัดแย้ง}
\code{scheduleProject} สร้างหน่วยสอบพร้อมกันจากความสัมพันธ์ของกฎ หรือใช้รายการสอบเดี่ยวเมื่อไม่เชื่อมกับวิชาอื่น น้ำหนักของหน่วยคือจำนวนรายการที่ต้องจัดและยังไม่มีเวลา ไม่ใช่จำนวนนักศึกษา จำนวนตอนเรียน หรือจำนวนแถวในแฟ้ม หน่วยที่มีรายการพร้อมเวลาอยู่แล้วถูกจำกัดให้ตรงกับเวลานั้น

ช่องเวลามาจากผลคูณของวันในช่วงสอบกับช่วงเวลารายวัน สำหรับการสอบเต็มวันใช้ช่วง fullDay จากนั้นกรองช่องที่ขัดกับข้อมูลคงที่ กลุ่มนักศึกษา ขอบเขตวัน และนโยบายปฏิทิน นโยบาย \code{never} ตัดวันออก ส่วนนโยบาย \code{only_if_necessary} ยังยอมให้ใช้ แต่คิดเป็นวันที่ควรหลีกเลี่ยง กราฟข้อขัดแย้งเชื่อมหน่วยที่มีกลุ่มนักศึกษาร่วมกันและมีช่องที่อาจทับซ้อน การแยกกราฟเป็นองค์ประกอบช่วยจำกัดการค้นหาเฉพาะส่วนที่เกี่ยวข้องกัน
ช่วงสอบรวมวันเริ่มและวันสิ้นสุด โดยไม่รับช่วงกลับด้านหรือยาวเกิน 366 วัน ช่วงเวลารายวันต้องถูกต้องอย่างน้อยหนึ่งช่วง แต่ไม่บังคับให้ต่างกันหรือไม่ทับซ้อน จึงตรวจข้อขัดแย้งจากช่วงนาทีจริง วันหยุดต้องระบุใน settings และไม่ได้ดึงจากปฏิทินราชการโดยอัตโนมัติ ค่าการตั้งค่าทั้งหมดอยู่ในภาคผนวก \ref{app:input}

\begin{figure}[tbp]\centering
\begin{tikzpicture}[x=0.023cm,y=1cm]
\draw[arr](500,0) to (990,0);
\foreach \x/\label in {540/09:00,720/12:00,780/13:00,960/16:00}{\draw(\x,0.1) to (\x,-0.1)node[below]{\label};}
\draw[line width=6pt,blue!55](540,0.65) to (720,0.65);
\draw[line width=6pt,orange!70](780,0.65) to (960,0.65);
\draw[line width=6pt,gray](540,1.35) to (960,1.35);
\node[above] at (750,1.4){เต็มวัน};
\end{tikzpicture}
\caption{ช่วงเต็มวันทับซ้อนกับทั้งช่วงเช้าและบ่ายในวันเดียวกัน}
\label{fig:overlap}
\end{figure}
\subsection{ตัวอย่างการค้นหาทีละขั้น}
ตัวอย่างสังเคราะห์ต่อไปนี้แสดงกลไกหลักในขนาดเล็กพอให้ติดตามการค้นหาได้ทีละขั้น กำหนดรายวิชาสี่วิชา สองวิชาถูกกฎผูกให้สอบพร้อมกันจึงรวมเป็นหน่วยเดียว มีสามหน่วย ช่วงสองสองช่วงในสามวัน และกำหนดให้วันที่สองเป็นวันที่ควรหลีกเลี่ยงตามนโยบายวันหยุด ลำดับการตัดสินใจของการค้นหาในตัวอย่างนี้เป็นดังนี้

\begin{figure}[tbp]\centering
\begin{tikzpicture}
\node[box,text width=3.6cm] (u1) at (0,0) {หน่วย U1 น้ำหนัก 2\\ รายวิชา A1 และ A2\\ กฎสอบพร้อมกัน\\ กลุ่มนักศึกษา G1};
\node[box,text width=3.6cm] (u2) at (6.4,0) {หน่วย U2 น้ำหนัก 1\\ รายวิชา B\\ กลุ่ม G1 และ G2};
\node[box,text width=3.6cm] (u3) at (3.2,-4.2) {หน่วย U3 น้ำหนัก 1\\ รายวิชา C เต็มวัน\\ กลุ่ม G2};
\draw[thick](u1) to node[above]{G1 ร่วม}(u2);
\draw[thick](u2) to node[above,sloped]{G2 ร่วม}(u3);
\end{tikzpicture}
\caption{กราฟข้อขัดแย้งของตัวอย่าง ขอบเกิดจากกลุ่มนักศึกษาร่วมเมื่อช่องเวลาของสองหน่วยทับซ้อนกันได้}
\label{fig:worked}
\end{figure}
\begin{center}\begin{tabular}{L{1.0cm}L{2.9cm}L{1.9cm}L{2.5cm}L{4.6cm}}\toprule
ขั้น & โดเมนที่เหลือ U1, U2, U3 & หน่วยที่เลือก & ช่องที่เลือก & เหตุผลของการเลือก \\\midrule
0 & 6, 6, 3 & U3 & วันที่ 1 เต็มวัน & MRV เลือกโดเมนน้อยที่สุด ต้นทุนเท่ากันจึงใช้วันแรก \\
1 & 6, 4, 0 & U2 & วันที่ 3 ช่วงเช้า & วันที่ 1 ถูกบล็อกโดย U3 วันที่ 2 คือวันที่ควรหลีกเลี่ยง จึงทำให้ต้นทุนองค์ประกอบแรกเพิ่มขึ้น และถูกเลื่อนไปท้ายลำดับ \\
2 & 5, 0, 0 & U1 & วันที่ 1 ช่วงเช้า & วันที่ 3 ช่วงเช้าถูกบล็อกโดย U2 ช่องที่เหลือที่ต้นทุนต่ำสุดคือวันแรก \\
หยุด & & & & น้ำหนักรวม 4 เท่ากับขอบเขตบน 4 การค้นหาจบโดยไม่ต้องสำรวจสาขาที่เหลือ \\\bottomrule
\end{tabular}\end{center}
ในทุกขั้น การเลือกหน่วยใช้ MRV เป็นเกณฑ์แรก และการเรียงช่องใช้ต้นทุนตามลำดับความสำคัญก่อนวันที่และเวลา การค้นหายังมีสาขาที่ไม่วางหน่วยเลย เพราะการไม่วางหน่วยเต็มวันหนึ่งอาจเปิดทางให้จัดรายการอื่นได้มากกว่า ตัวอย่างนี้สังเคราะห์ขึ้นเพื่อการอธิบาย ชุดข้อมูลหลักที่ใช้ทดสอบในรายงานมีขนาดและความเชื่อมโยงมากกว่าตัวอย่างนี้มาก ดังสถิติโครงสร้างในบทที่ \ref{ch:method}
\subsection{การเลือกหน่วยถัดไปและการเรียงช่องเวลา}
\code{solveUnits} ใช้ MRV เป็นเกณฑ์แรก โดยนับช่องที่ยังไม่ถูกบล็อกในสถานะปัจจุบัน เมื่อจำนวนเท่ากัน จึงเรียงน้ำหนักมากก่อน สอบเต็มวันก่อน จำนวนเพื่อนบ้านในกราฟมากก่อน จำนวนสมาชิกกลุ่มมากก่อน และดัชนีคงที่น้อยก่อน หน่วยถูกเรียงตาม ID ก่อนสร้างดัชนี จึงมีตัวตัดสินสุดท้ายที่ทำซ้ำได้

เมื่อความครอบคลุมเท่ากัน ระบบลดต้นทุนแบบเรียงลำดับความสำคัญสามองค์ประกอบ ได้แก่ จำนวนรายการในวันที่ควรหลีกเลี่ยง ภาระของกลุ่มนักศึกษา และความต้องการใช้ทรัพยากรพร้อมกันในช่วงเวลาเดียวกัน การเปรียบเทียบไม่ใช่การรวมทั้งสามเป็นคะแนนเดียว จึงไม่ยอมเพิ่มองค์ประกอบแรกเพียงเพื่อให้องค์ประกอบท้ายลดลงมาก กรอบปัญหา CSP และแนวคิด MRV ที่ใช้ในบทนี้อ้างอิงจาก Russell และ Norvig (2021) และโน้ตประกอบการสอนของ Carnegie Mellon University (n.d.) ส่วนต้นทุนแบบเรียงลำดับความสำคัญ สัมประสิทธิ์ และขอบเขตการพิสูจน์ใช้ implementation ของ \code{solveUnits} เป็นนิยามเฉพาะระบบ ไม่อ้างว่าตำรารับรองอัลกอริทึมนี้ทั้งหมด

ภาระที่เพิ่มจากการเลือกช่องคำนวณจากสมาชิกกลุ่มของรายการที่กำลังจัด โดยให้น้ำหนัก 10 ต่อรายการร่วมกลุ่มในวันเดียวกัน และ 3 ต่อรายการในวันปฏิทินก่อนหน้าหรือถัดไป น้ำหนักสองค่านี้เป็นพารามิเตอร์ที่กำหนดพฤติกรรม ไม่ใช่ค่าที่พิสูจน์ทางคณิตศาสตร์ เหตุผลของอัตราส่วนคือการที่การสอบร่วมวันเดียวกันกดดันนักศึกษามากกว่ากรณีคนละวันที่ติดกันหลายเท่า ผลของการเปลี่ยนค่าเหล่านี้วัดไว้ในบทที่ \ref{ch:perf} แล้วพบว่าอัตราส่วนที่ใกล้เคียงกันให้ผลแทบไม่ต่างกัน ส่วนองค์ประกอบที่สามนับจำนวนรายการ managed ที่เวลาทับซ้อนคูณน้ำหนักหน่วย เป็นตัวแทนของความต้องการใช้ห้องพร้อมช่วงเวลากันเท่านั้น ไม่ใช่การจัดห้องสอบหรือการพิจารณาความจุห้องจริง การลดค่านี้ช่วยลดโอกาสที่ตารางจะต้องใช้ห้องพร้อมกันจำนวนมากในช่วงเดียว เมื่อค่าต้นทุนเพิ่มเท่ากัน เรียงวันที่ เวลาเริ่ม และเวลาสิ้นสุดจากน้อยไปมาก
\subsection{งบการค้นหา การตัดกิ่ง และการปรับปรุง}
\needspace{10\baselineskip}แผนการค้นหาโดยสรุปมีดังนี้

สำหรับแต่ละองค์ประกอบของกราฟข้อขัดแย้ง: คำนวณขอบเขตบนของความครอบคลุม จากการจัดเรียงช่วงเวลาของแต่ละกลุ่มนักศึกษา ค้นหาความเป็นไปได้โดยเลือกหน่วยด้วย MRV ทีละตัว: ลองช่องเวลาตามต้นทุนจากน้อยไปมาก โดยเลือกวันปกติก่อนวันที่ควรหลีกเลี่ยง งดสำรวจสาขาที่ศักยภาพรวมต่ำกว่าคำตอบที่ดีที่สุดปัจจุบัน (ตัดกิ่งด้วยขอบเขตบน) พิจารณาสาขาที่ไม่วางหน่วยเมื่อการนั้นอาจเพิ่มความครอบคลุมรวม เมื่อสิ้นสุดการค้นหาความเป็นไปได้ ทดลองคำตอบตั้งต้นแบบ greedy อีกสามลำดับที่กำหนดแน่นอน ปรับปรุงภายในงบที่เหลือ ด้วยการซ่อมความครอบคลุม การย้ายหน่วย และการแลกที่อยู่สองหน่วย เก็บคำตอบที่ดีที่สุดตามน้ำหนักก่อน แล้วจึงต้นทุนตามลำดับความสำคัญ

งบค้นหาความเป็นไปได้เริ่มที่ค่ามากกว่าระหว่าง 1 กับจำนวนเต็มปัดลงของร้อยละ 60 ของ \code{searchBudget} และแบ่งระหว่างองค์ประกอบตามจำนวนตัวแปร สัดส่วนนี้เป็นพารามิเตอร์เช่นกัน การเปลี่ยนสัดส่วนไป 0.4 หรือ 0.8 ได้ผลชนิดเดียวกันตามที่วัดไว้ในบทที่ \ref{ch:perf} โดยเปลี่ยนเพียงการแบ่งงบระหว่างการค้นหากับการปรับปรุง ระบบพิจารณาวันที่ควรใช้ก่อน แล้วขยายไปยังวันที่ควรหลีกเลี่ยงเมื่อจำเป็น

การตัดกิ่งด้วยขอบเขตบนใช้ค่าที่ประเมินอย่างผ่อนปรนจากช่องเวลาและกลุ่มนักศึกษา หากสาขานั้นไม่อาจเพิ่มความครอบคลุมหรือปรับจำนวนรายการในวันที่หลีกเลี่ยงก็ไม่จำเป็นต้องสำรวจต่อ องค์ประกอบที่ถูกงบตัดไปแล้วจะไม่ถูกซ่อมให้ดูสมบูรณ์ในเบื้องหลัง

\code{visitedNodes} นับงานค้นหา ส่วน \code{improvementEvaluations} นับการประเมินปรับปรุง ผลรวมไม่เกินงบ แต่ไม่ครอบคลุมทุกคำสั่งประมวลผล เช่น การสร้างคำตอบตั้งต้นบางส่วน จึงไม่ควรแปลงตัวนับเหล่านี้เป็นเวลาหรือจำนวนการทำงานทั้งหมด
\subsection{การคำนวณขอบเขตบนของความครอบคลุม}
ขอบเขตบนคำนวณแยกตามองค์ประกอบของกราฟข้อขัดแย้ง เริ่มจากน้ำหนักรวมของทุกหน่วยในองค์ประกอบ แล้วหักส่วนที่ต้องเสียไปตามข้อจำกัดของแต่ละกลุ่มนักศึกษา สำหรับกลุ่มหนึ่ง รวบรวมช่องเวลาที่สมาชิกของกลุ่มใช้ได้ แล้วนับจำนวนช่องที่ไม่ทับซ้อนกันตามลำดับวันและเวลาสิ้นสุด จำนวนนี้คือความจุสูงสุดของกลุ่มนั้น หากกลุ่มมีหน่วยมากกว่าความจุ ส่วนเกินต้องหลุดออก โดยตัดน้ำหนักน้อยที่สุดออกก่อน ผลต่างที่เหลือคือความสูญเสียขั้นต่ำของกลุ่มนั้น ขอบเขตขององค์ประกอบจึงเท่ากับน้ำหนักรวมลบความสูญเสีย โดยการรวมค่าความสูญเสียของหลายกลุ่มเลือกเฉพาะกลุ่มนักศึกษาที่ไม่มีหน่วยสอบร่วมกัน เพื่อไม่ให้รายการสอบเดียวกันถูกนับเป็นความสูญเสียมากกว่าหนึ่งครั้ง

ขอบเขตนี้ไม่เคยต่ำกว่าคำตอบจริง เพราะทุกตารางที่เป็นไปได้ต้องวางสมาชิกของกลุ่มภายในช่องที่ใช้ได้ และช่องที่ทับซ้อนกันรองรับได้เพียงหนึ่งหน่วย จึงวางได้ไม่เกินความจุที่คำนวณแบบมองโลกในแง่ดีนี้ เมื่อคำตอบที่ค้นพบมีน้ำหนักเท่ากับขอบเขตบนของทุกองค์ประกอบรวมกัน จึงไม่มีคำตอบใดดีกว่าได้ และ \code{coverageOptimal} เป็นจริง สำหรับจำนวนรายการในวันที่ควรหลีกเลี่ยง ขอบเขตล่างคือผลต่างระหว่างขอบเขตบนรวมช่องทั้งหมดกับขอบเขตบนเฉพาะวันปกติ หากคำตอบที่ครอบคลุมถึงขอบเขตบนแล้วยังมีจำนวนรายการในวันดังกล่าวเท่ากับขอบเขตล่างนี้ \code{avoidedDaysOptimal} จึงเป็นจริงตามนิยามใน \code{src/lib/solver.ts}
\needspace{8\baselineskip}\subsection{ลำดับค้นหาและขอบเขตของข้อพิสูจน์}
\code{coverageOptimal} หมายถึงจำนวนที่จัดได้ถึงขอบเขตที่ใช้ได้ หรือค้นหาครบจนพิสูจน์ความครอบคลุมได้ ไม่ได้ยืนยันว่าภาระนักศึกษาต่ำที่สุดทุกด้าน \code{avoidedDaysOptimal} เพิ่มเงื่อนไขว่าจำนวนรายการในวันที่ควรหลีกเลี่ยงถึงขอบเขตล่างด้วย

\noindent\code{budgetExhausted} ระบุการหยุดค้นหาความเป็นไปได้โดยยังไม่มีข้อพิสูจน์ความครอบคลุม เหตุการณ์เช่นนี้หมายความว่าไม่พบคำตอบภายในงบที่ตั้งไว้เท่านั้น ไม่ได้แปลว่าปัญหาไม่มีคำตอบ และไม่ใช่หลักฐานว่าจัดต่อไม่ได้แม้เพิ่มงบ

\code{provablyImpossibleEvents} ระบุรายการเฉพาะที่จัดไม่ได้ภายใต้ข้อจำกัดที่พิสูจน์ได้ รายการเหล่านี้เป็นข้อเท็จจริงเกี่ยวกับข้อจำกัด ไม่ใช่ข้อสรุปว่าโครงการทั้งหมดไม่มีตารางที่เป็นไปได้ หากข้อมูลนำเข้าขัดกัน ส่วนที่ขัดกันอาจคงอยู่ในผลลัพธ์และทำให้ผลตรวจรวมเป็น invalid
\begin{figure}[tbp]\centering
\begin{tikzpicture}[node distance=0.5cm,box/.append style={text width=3.5cm,minimum height=1.2cm}]
\node[box] (a) {สร้างหน่วย ช่องเวลา และกราฟ};
\node[box,right=of a] (b) {เลือกหน่วยด้วย MRV};
\node[box,right=of b] (c) {ลองช่องตามต้นทุนและลำดับคงที่};
\node[box,below=of c] (d) {ปรับขอบเขตและตัดกิ่ง};
\node[box,left=of d] (e) {เก็บคำตอบและปรับปรุงภายในงบ};
\node[box,left=of e] (f) {รายงานผลพร้อมสถานะข้อพิสูจน์};
\foreach \x/\y in {a/b,b/c,c/d,d/e,e/f}\draw[arr](\x) to (\y);
\end{tikzpicture}
\caption{การค้นหาสร้างคำตอบและรายงานขอบเขตความมั่นใจอย่างไร}
\label{fig:solverflow}
\end{figure}

\chapter{การตรวจสอบความถูกต้อง}\label{ch:valid}
\subsection{การตรวจระหว่างสร้างกับการตรวจซ้ำ}
ระหว่างค้นหา ระบบกรองช่องที่ขัดกับข้อจำกัดและป้องกันการชนของกลุ่ม แต่การตรวจหลังได้ผลยังจำเป็น เพราะผลอาจมีเวลาที่นำเข้าอยู่ก่อน มีการแก้ด้วยมือ หรือมีระเบียนที่ไม่ตรงกับต้นทาง \code{validateProject} จึงสร้างรายการมาตรฐานใหม่ผ่าน \code{buildEvents} และใช้รายการดังกล่าวเป็นฐานตรวจผลจริง

การตรวจซ้ำผ่านเส้นทางที่แยกจากตัวค้นหา และการตรวจจับผลลัพธ์ที่ถูกแก้ไข ทำให้ตรวจพบรายการหาย รายการเกิน ID ซ้ำ และเขตข้อมูลที่ไม่ตรงกับต้นทางได้ การแก้กลุ่มนักศึกษาในผลลัพธ์เพียงอย่างเดียวไม่ควรซ่อนการชน เพราะตัวตรวจใช้กลุ่มที่คำนวณใหม่ อย่างไรก็ตาม กลไกนี้ยังเชื่อถือข้อมูลต้นทาง ไม่ใช่ระบบพิสูจน์ความแท้ของระเบียนจากภายนอก
\subsection{กลุ่มการตรวจที่อยู่ในขอบเขต}
\begin{longtable}{L{4.4cm}L{8.6cm}}\toprule
ชื่อกลุ่ม & ประเด็นที่ตรวจ \\\midrule
\endfirsthead
\toprule
ชื่อกลุ่ม & ประเด็นที่ตรวจ \\\midrule
\endhead
\code{imports} & ปัญหาจากการนำเข้าและการประสานกฎ \\
\code{canonicalEvents} & ความครบและความตรงกับรายการที่สร้างจากต้นทาง \\
\code{studentOverlap} & การสอบทับซ้อนของกลุ่มนักศึกษา \\
\code{sameTimeRules} & รายการตามกฎมีวันเวลาเดียวกัน \\
\code{dateRanges} & วันสอบอยู่ภายในช่วงที่กำหนด \\
\code{requiredExams} & รายการที่ต้องสอบมีวันเวลา \\
\code{fixedTimings} & เวลาเดิมและการล็อกถูกเก็บรักษา \\
\code{fullDay} & ช่วงเวลาตรงกับการสอบเต็มวัน \\
\code{calendarPolicy} & วันหยุดและวันสุดสัปดาห์ตามนโยบาย \\
\code{resourceIntegrity} & รหัสและการอ้างอิงถูกต้อง \\\bottomrule
\end{longtable}
ตารางนี้มี 10 กลุ่ม ภายในขอบเขตรายงาน ไม่ใช่จำนวนเงื่อนไขย่อยหรือจำนวนกรณีทดสอบ หน้าจอปัจจุบันกรองบางกลุ่มออก จึงไม่ควรเปรียบเทียบจำนวนแถวที่มองเห็นกับตารางนี้โดยตรง สถานะ \code{not_checked} ของกลุ่มใดหมายความว่าไม่มีข้อมูลเพียงพอจะตรวจ ไม่ใช่ผ่านหรือไม่ผ่าน

\needspace{24\baselineskip}\subsection{ระดับความรุนแรงและสถานะ}
ข้อค้นพบมีระดับ \code{error} \code{warning} และ \code{info} พร้อมที่มา \code{source} \code{generated} \code{manual} หรือ \code{resource} สถานะรวมเป็น \code{invalid} เมื่อมี \code{error} อย่างน้อยหนึ่งรายการ หากไม่มี \code{error} แต่ยังมีข้อค้นพบ จะเป็น \code{valid_with_warnings} แม้ข้อค้นพบนั้นจะเป็น \code{info} เท่านั้น จึงควรอ่านสถานะร่วมกับรายการจริง

แต่ละกลุ่มตรวจมีสถานะ \code{passed} \code{failed} หรือ \code{not_checked} การพบข้อค้นพบที่อยู่ในกลุ่มทำให้กลุ่มนั้น \code{failed} ได้แม้เป็นคำเตือน ส่วนกลุ่มที่ไม่มีข้อมูลเพียงพอและไม่มีข้อค้นพบอาจเป็น \code{not_checked}
\needspace{24\baselineskip}\subsection{ข้อผิดพลาดที่มีรหัสเหมือนกันต่างบทบาท}
\code{SAME_TIME_CONTRADICTION} ปรากฏได้ทั้งในฐานะ issue ของตัวตรวจ และเหตุผลของหน่วยที่ตัวจัดตารางวางไม่ได้ การจำแนกจึงใช้บริบทของเขตข้อมูลร่วมกับตัวอักษรรหัส ไม่กำหนดว่ารหัสหนึ่งอยู่ในหมวดเดียวเสมอ เช่น \code{issueCodes} และ \code{solverUnscheduledReasons} อาจมีข้อความเดียวกัน แต่บอกข้อเท็จจริงคนละระดับ

ในชุดขัดแย้ง วิชาคู่เดียวมีกลางภาคและปลายภาค จึงเกิดเหตุผลสองหน่วยแต่มีรายการสอบไม่มีเวลาสี่รายการ หากรายงานนับสองเป็นจำนวนรายการสอบจะทำให้ขนาดผลกระทบลดลงครึ่งหนึ่ง เครื่องมือยืนยันผลจึงตรวจทั้งสองจำนวนแยกกัน
\needspace{20\baselineskip}\subsection{แนวทางใช้ผลตรวจ}
ผลตรวจควรนำไปแก้ต้นเหตุ ปัญหารูปแบบข้อมูลแก้ในแฟ้มต้นทาง ปัญหากฎแก้รายการอ้างอิงหรือความหมายของกฎ ส่วนงบค้นหาหมดต้องพิจารณาทั้งโครงสร้างข้อมูลและงบที่อนุญาต การเพิ่มงบอาจช่วยค้นหาต่อ แต่ไม่แก้รหัสวิชาผิดรูปแบบหรือข้อกำหนดที่ขัดกัน

รายงานไม่ถือว่าการผ่าน test suite ยืนยันทุกชุดข้อมูลจริง การทดสอบโปรแกรมเดิมกับการตรวจ fixture เป็นหลักฐานคนละบทบาท ดังที่แยกไว้ในบทที่ \ref{ch:method} ทั้งสองส่วนจึงต้องอ่านร่วมกันเมื่ออ้างผล
\chapter{วิธีการทดสอบและชุดข้อมูลสังเคราะห์}\label{ch:method}
\subsection{ช่องทางหลักฐานและรหัสออก}
การเรียกจัดตารางแบบเต็มพิมพ์ผลตรวจต้นทางไปยัง \code{stderr} และพิมพ์สรุปหลังทำงานไปยัง \code{stdout} ส่วนโหมดตรวจนำเข้าพิมพ์ผลตรวจไปยัง \code{stdout} เครื่องมือทดสอบจึงเก็บสองช่องทางแยกกัน การแยกนี้สำคัญต่อการวินิจฉัยระยะที่เกิดปัญหา หากมีเพียงข้อความผิดพลาดจากการอ่านแฟ้ม ย่อมอ้างไม่ได้ว่าตัวแก้ปัญหาเริ่มค้นหาแล้ว หากมีผลลัพธ์เต็มพร้อมรายงานตรวจซ้ำ ก็ยืนยันได้ว่าระบบไปถึงขั้นสร้างหลักฐาน แม้รหัสออกจะเป็น 2 ก็ตาม การมีแฟ้มผลลัพธ์ไม่ได้หมายถึงผลผ่านการตรวจ เพราะ CLI เขียนแฟ้มก่อนกำหนดสถานะออกจากโปรเซสในกรณีที่ตารางยังไม่สมบูรณ์

\begin{center}\begin{tabular}{L{2.2cm}L{10.6cm}}\toprule
รหัสออก & ความหมายในกระบวนการทดลอง \\\midrule
0 & ทำงานสำเร็จและผลไม่อยู่ในสถานะ invalid อาจยังมีคำเตือน \\
1 & ความล้มเหลวด้านการเรียกใช้ การอ่านแฟ้ม หรือการประมวลผลที่โยนข้อผิดพลาด \\
2 & พบปัญหาข้อมูลที่บังคับให้หยุด หรือได้ผลตารางที่ไม่ครบหรือไม่ถูกต้อง \\\bottomrule
\end{tabular}\end{center}
ค่าแฮชของตารางช่วยตรวจความสม่ำเสมอของผลที่ได้ แต่ไม่ใช่การพิสูจน์ความถูกต้องของข้อกำหนดทั้งหมด การทำซ้ำจึงเปรียบเทียบทั้งแฮช รายการที่จัดไม่ได้ และสถิติการค้นหา ควบคู่กับเงื่อนไขเชิงพฤติกรรม เช่น วิชาที่ต้องสอบพร้อมกันมีวันเวลาเท่ากันจริง
\subsection{การตรวจยืนยันผล}
\code{verify-results.ts} ตรวจรหัสออก จำนวนข้อค้นพบ สถานะ จำนวนรายการที่จัดได้ เหตุผลรายการที่ขาด และเงื่อนไขหลัก จากนั้นเรียกชุด clean ซ้ำ และเปรียบเทียบ schedule hash รายการที่จัดไม่ได้ และสถิติการค้นหา ค่าใน \code{expected.json} มาจากผลที่สังเกตก่อนตรึงความคาดหมาย การตรวจครั้งใหม่ไม่ปรับค่าคาดหมายเพื่อให้ผ่านโดยไม่มีเหตุผล การเทียบกับค่าที่ตรึงไว้จึงเป็นหลักฐานความเสถียรของการทำซ้ำ (regression test) ไม่ใช่หลักฐานอิสระว่าค่าที่ตรึงไว้ครั้งแรกถูกต้อง ความถูกต้องของผลอ้างชั้นตรวจเชิงพฤติกรรมและเงื่อนไขความถูกต้องที่ตรวจแยกออกจากค่าที่ตรึงไว้ ทั้งสองชั้นหลักฐานนี้ต่างบทบาทกัน การทดสอบชุดนี้รันซ้ำที่ revision ปัจจุบันก่อนสร้างรายงานฉบับนี้และผ่านทั้งแปดสถานการณ์

\subsection{โครงสร้างตัวสร้างข้อมูล}
\code{docs/data/generate.ts} สร้าง workbook รายวิชาและ CSV สำหรับกฎกับห้อง โดยใช้ค่าตั้งต้นและลำดับสุ่มที่กำหนดแน่นอน รายชื่อและข้อมูลตัวอย่างเป็นข้อมูลสังเคราะห์ แฟ้มรายงานแยกจากตัวอย่างในผลิตภัณฑ์ซึ่งสร้างโดย \code{scripts/make-example-fixtures.ts} จึงไม่เขียนทับ fixture ที่โปรแกรมใช้ทดสอบอยู่เดิม

ชุดหลักมี managed 41 วิชา 47 ตอนเรียน และ context 10 วิชา ขนาดกลางมี managed 100 วิชา 126 ตอนเรียน กับ context 59 วิชา 85 ตอนเรียน ขนาดใหญ่มี managed 189 วิชา 238 ตอนเรียน กับ context 111 วิชา 160 ตอนเรียน จำนวนวิชารวมจึงเป็น 51, 159 และ 300 ตามลำดับ การระบุ managed ควบคู่กับยอดรวมช่วยป้องกันการเปรียบเทียบคนละหน่วย

\subsection{การกำหนดขนาดชุดข้อมูล}
ขนาดของชุดข้อมูลสังเคราะห์ถูกเลือกให้ใกล้เคียงกับขนาดงานที่ระบบมีเป้าหมายรองรับ โดยอ้างบริบทของขนาดจากการทดลองนำร่องกับข้อมูลต้นทางจริงของหน่วยงานต้นแบบ ข้อมูลจริงไม่ได้ถูกใช้เป็นผลการทดลองของรายงาน และไม่เผยแพร่ผลรายบุคคลในโฟลเดอร์รายงาน ผลการทดลองทุกตัวเลขในรายงานฉบับนี้จึงมาจากชุดข้อมูลสังเคราะห์เท่านั้น

หน้าจอรายวิชาแสดงรายการสอบที่ระบบรับผิดชอบ พร้อมวันเวลา กลุ่มนักศึกษา และป้ายบทบาทที่มาของแต่ละรายการ
\begin{figure}[tbp]\centering
\includegraphics[width=\textwidth]{docs/screenshots/exams.png}
\caption[หน้าจอรายวิชาจากชุดข้อมูลหลัก]{หน้าจอรายวิชาจากชุดข้อมูลหลัก แสดงรายการสอบที่จัดแล้วพร้อมวันเวลาและกลุ่มนักศึกษา}
\label{fig:exams}
\end{figure}
\FloatBarrier
\subsection{โครงสร้างของชุดข้อมูลเทียบกับความยากของปัญหา}
จำนวนวิชาไม่ใช่ตัวกำหนดความยากของปัญหาโดยลำพัง ตารางนี้จำแนกชุดข้อมูลด้วยโครงสร้างข้อกำหนดที่ตัวแก้ปัญหาเผชิญจริง คำนวณโดยสคริปต์ \code{docs/data/fixture-stats.ts} ซึ่งทำซ้ำการสร้างหน่วยและกราฟข้อขัดแย้งจากซอร์สโค้ดตัวจัดตาราง และตรวจยืนยันความสอดคล้องของสำเนาโดยเทียบกับผลลัพธ์ที่บันทึกไว้ก่อนเขียนตัวเลขใด ๆ

\begin{center}\begin{minipage}{\textwidth}\centering
\captionof{table}{สถิติโครงสร้างของชุดข้อมูล}\label{tab:structure}
\begin{tabular}{L{5.4cm}ccc}\toprule
ตัวชี้วัด & ชุดเล็ก & ชุดกลาง & ชุดใหญ่ \\\midrule
รายการสอบ managed & 82 & 200 & 378 \\
มีวันเวลาจากต้นทางแล้ว & 4 & 0 & 0 \\
รายการสอบเต็มวัน & 2 & 4 & 4 \\
รายการ context ที่มีเวลา & 20 & 118 & 222 \\
หน่วยค้นหารวม (รวมกฎสอบพร้อมกัน) & 76 & 194 & 372 \\
กลุ่มนักศึกษาที่ต้องประสาน & 20 & 24 & 24 \\
ขอบข้อขัดแย้งในกราฟ & 300 & 984 & 3338 \\
ความหนาแน่นของกราฟ (สัดส่วนขอบที่เป็นไปได้) & 10.5\% & 5.3\% & 4.8\% \\
องค์ประกอบของกราฟ & 16 & 2 & 2 \\
ระดับการเชื่อมโยงเฉลี่ยของหน่วย & 7.9 & 10.1 & 17.9 \\
ช่องเวลาที่แตกต่างกันทั้งหมด & 75 & 87 & 74 \\
ห้องทั้งหมด (ติดป้ายคอมพิวเตอร์) & 16 (3) & 16 (3) & 24 (4) \\
งบการค้นหา & 20000 & 20000 & 20000 \\\bottomrule
\end{tabular}\end{minipage}\end{center}
จำนวนวันในช่วงสอบเท่ากันทุกชุด คือ 12 วันกลางภาคและ 17 วันปลายภาคตามปฏิทิน ตารางนี้อธิบายเหตุผลหลักของเวลาที่วัดได้ในบทที่ \ref{ch:perf} ได้ดีกว่าจำนวนวิชาเพียงตัวเดียว กล่าวคือ ชุดที่ใหญ่ขึ้นมีระดับการเชื่อมโยงต่อหน่วยสูงขึ้นมาก ขณะที่ความหนาแน่นโดยรวมของกราฟลดลง และองค์ประกอบที่แตกกระจาย 16 ก้อนของชุดเล็กรวมตัวเป็นกราฟใหญ่เกือบเดี่ยวในชุดกลางและใหญ่

ชุดใหญ่ใช้ seed 5 ซึ่งเลือกเป็น seed แรกที่ทุกกลุ่มนักศึกษามีความต้องการไม่เกินจำนวนช่องเวลาในช่วงสอบ โดยรักษาขนาดเดิมและไม่ได้เลือกจากความเร็วของผลวัด ชุดวัดเวลาจึงมีเงื่อนไขอยู่บนความสามารถจัดครบโดยโครงสร้าง ไม่ใช่ตัวอย่างแบบไม่มีอคติจากการกระจายความยากทั้งหมดของตัวสร้างข้อมูล ตัวเลขเวลาในบทที่ \ref{ch:perf} จึงอ่านได้ในความหมายของชุดที่จัดครบได้เท่านั้น
\needspace{8\baselineskip}\subsection{กรณีหลักที่ตรวจเชิงพฤติกรรม}
\begin{center}\begin{minipage}{\textwidth}\centering
\captionof{table}{กรณีทดสอบเชิงพฤติกรรม}\label{tab:behavior}
\begin{tabular}{L{3.2cm}L{4.4cm}L{5.6cm}}\toprule
โครงสร้าง & คุณสมบัติที่ตรวจ & ผลคาดหมายและผลที่พบ \\\midrule
วิชา \code{060321113} กับ \code{060321120} & กฎสอบพร้อมกัน & วันและช่วงตรงกันทั้งสองประเภทสอบ \\
วิชา \code{060321127} กับ \code{060321134} & กลุ่มนักศึกษาร่วมกัน & ไม่มีช่วงทับซ้อน \\
วิชา \code{060321106} & สอบเต็มวัน & เริ่ม 540 สิ้นสุด 960 นาที \\
กลุ่มชั้นปีที่ 2 ร่วมสี่กลุ่ม & วันธรรมดาไม่พอ & รายการในวันที่หลีกเลี่ยงอยู่ในกลุ่มที่อิ่มตัว \\
วิชา \code{060321190} กับ \code{060321197} & รักษาเวลาต้นทาง & วันเวลาเท่าเดิมและมีที่มา imported \\
วิชาหกตอนเรียน & การรวมระเบียน & ไม่สร้างรายการสอบแยกตามทุกตอนเรียน \\\bottomrule
\end{tabular}\end{minipage}\end{center}
กลุ่มอิ่มตัวในที่นี้หมายถึงกลุ่มนักศึกษาที่มีรายการสอบรวมมากกว่าจำนวนช่วงเวลาวันปกติที่ยังว่างอยู่ กลุ่มนี้ต้องการช่วงเวลารวม 17 หน่วยช่วงเวลา แต่รายการประกอบยึดช่วงเช้าวันธรรมดาไว้จนเหลือ 12 หน่วย การใช้วันที่ควรหลีกเลี่ยงจึงเกิดจากความต้องการของกลุ่มเอง ไม่ใช่เพียงผลจากการเลือกวันที่ไม่ดี ตัวตรวจยืนยันว่ารายการในวันดังกล่าวเกี่ยวข้องกับกลุ่มนี้จริง
\needspace{10\baselineskip}\subsection{ชุดข้อผิดพลาดและระยะตรวจพบ}
\begin{center}\begin{minipage}{\textwidth}\centering
\captionof{table}{ข้อผิดพลาดและระยะตรวจพบ}\label{tab:errors}
\begin{tabular}{L{3.0cm}L{3.0cm}L{7.2cm}}\toprule
สาเหตุ & ระยะที่พบ & รหัสและผลกระทบ \\\midrule
อ้างวิชาที่ไม่มี & ประสานกฎ & \code{STALE_RULE_REFERENCE} เป็นคำเตือน ยังจัดได้ \\
วิชาที่ต้องตรงกันมีกลุ่มร่วม & จัดตารางและตรวจซ้ำ & \code{SAME_TIME_CONTRADICTION} ทำให้สี่รายการไม่มีเวลา \\
กฎซ้ำในชุดรวม & ประสานกฎ & \code{DUPLICATE_RULE} เป็นข้อมูลแจ้งให้ทราบ \\
รหัสวิชาไม่ครบเก้าหลัก & นำเข้ารายวิชา & \code{MALFORMED_COURSE_ROW} บล็อกการตรวจนำเข้า \\
ไม่มีกลุ่มนักศึกษา & นำเข้าและสร้างรายการ & \code{MISSING_STUDENT_GROUPS} และ \code{SOURCE_DATA_INVALID} \\
วันที่หรือช่วงเวลาผิด & นำเข้ารายวิชา & \code{INVALID_EXAM_TIMING} และ \code{SOURCE_DATA_INVALID} \\
ชื่อห้องซ้ำ & นำเข้าห้อง & \code{DUPLICATE_ROOM} บล็อกการตรวจนำเข้า \\\bottomrule
\end{tabular}\end{minipage}\end{center}
รายการ \code{INVALID_EXAM_TIMING} และ \code{MALFORMED_COURSE_ROW} ตรวจพบโดยตัวนำเข้ารายวิชาโดยตรง ก่อนเข้าสู่กระบวนการจัดตาราง ส่วนชุดสอบพร้อมกันขัดแย้งนำเข้าได้ แต่ล้มเหลวในความหมายของข้อกำหนดการจัดตาราง

หน้าจอข้อควรตรวจสอบแสดงข้อค้นพบพร้อมที่มาและระดับความรุนแรง สำหรับชุดกฎอ้างอิงล้าสมัย กฎที่อ้างรหัส \code{099999999} ซึ่งไม่มีในข้อมูลรายวิชาจะปรากฏเป็นคำเตือนที่มาของรายการคือแฟ้มต้นทาง
\begin{figure}[tbp]\centering
\includegraphics[width=\textwidth]{docs/screenshots/issues-stale-rule.png}
\caption[หน้าจอข้อควรตรวจสอบ ชุดกฎอ้างอิงล้าสมัย]{หน้าจอข้อควรตรวจสอบจากชุดกฎอ้างอิงล้าสมัย แสดงคำเตือนที่มีที่มาจากต้นทาง}
\label{fig:issues}
\end{figure}
\subsection{การปฏิเสธข้อมูลกับการโยนข้อผิดพลาด}
สี่ชุดปฏิเสธข้อมูลในงานนี้เรียกโหมด inspect และได้รหัสออก 2 ทั้งหมด การอ่านแฟ้มที่ว่าง ชนิดไม่รองรับ หรือใหญ่เกิน 25 MiB มีเส้นทางโยนข้อผิดพลาดและออกด้วย 1 ตามซอร์สโค้ด แต่ยังไม่ได้ใช้เป็นหนึ่งในสี่ fixture นี้ จึงแยกคำอธิบายความสามารถจากผลที่วัดจริง
\subsection{ภาพหน้าจอเป็นหลักฐาน}
ภาพหน้าจอทั้งแปดภาพบันทึกจากสภาพแวดล้อมเดียวกัน คือเบราว์เซอร์ที่ล้างข้อมูลการจัดเก็บ ภาษาไทยเป็นค่าเริ่มต้น ขนาดหน้าจอ 1920x1080 และระดับซูม 100 เปอร์เซ็นต์ รายละเอียดการบันทึกทั้งหมด เช่น รุ่นเบราว์เซอร์ เวลา และแฮช SHA-256 ของแฟ้มภาพ อยู่ในแฟ้ม \code{docs/screenshots/screenshots.json} และสรุปไว้ในภาคผนวก C ภาพปกติทั้งหมดใช้ชุดข้อมูลหลัก ส่วนหน้าจอข้อควรตรวจสอบใช้ชุดกฎอ้างอิงล้าสมัยตามที่ระบุในคำบรรยายของภาพ ดังที่อธิบายไว้ในบทที่ 1 ภาพเหล่านี้บันทึกจากจุดอ้างอิงของวันที่ 5 กันยายน 2569 ซึ่งโค้ดการจัดตารางไม่ต่างจาก revision ปัจจุบัน การโต้ตอบที่ลึกกว่าการเปิดดูยังไม่ถูกทดสอบและอยู่ในงานต่อไป ส่วนชั้น cloud ตรวจบนระบบที่เผยแพร่จริงด้วยหลักฐานของตัวเองในหัวข้อถัดไป
\needspace{8\baselineskip}\subsection{การตรวจสอบชั้น cloud บนระบบที่เผยแพร่จริง}\label{sec:cloud}
ชั้น cloud ตรวจบนระบบที่เผยแพร่จริงที่ \code{https://schedule.breadtm.xyz/} เมื่อวันที่ 6 กันยายน 2569 ด้วยเครื่องมือ \cmd{node docs/data/cloud-verification.ts} ซึ่งเรียก API ของ worker โดยตรง ข้อมูลที่ใช้คือโครงการมาตรฐานจากแฟ้ม \code{docs/data/results-clean.json} ก่อนเริ่มตรวจ เครื่องมือสร้างแฟ้มส่วนติดต่อผู้ใช้จาก revision ปัจจุบันแล้วเทียบกับแฟ้ม JavaScript CSS และนโยบายความปลอดภัยของเนื้อหาที่เซิร์ฟเวอร์จริงตอบกลับ ผลตรงกันทุกรายการ ผลการตรวจทั้งหมดบันทึกในแฟ้ม \code{docs/data/cloud-verification.json} และสรุปในตารางที่ \ref{tab:cloudchecks} ซึ่งผ่านเกณฑ์ครบทั้ง 16 รายการ

{\small\renewcommand{\arraystretch}{1.2}
\begin{longtable}{L{5.8cm}L{7.4cm}}
\caption[ผลการตรวจชั้น cloud]{ผลการตรวจชั้น cloud บนระบบที่เผยแพร่จริง สรุปจากแฟ้ม \code{docs/data/cloud-verification.json} ทั้ง 16 รายการผ่านเกณฑ์}\label{tab:cloudchecks}\\
\toprule รายการตรวจ & ผลที่สังเกต \\\midrule
\endfirsthead
\toprule รายการตรวจ & ผลที่สังเกต \\\midrule
\endhead
ความเหมือนของแฟ้มส่วนหน้าจอกับผลสร้างในเครื่อง & แฟ้ม JavaScript CSS และนโยบายความปลอดภัยของเนื้อหาตรงกันทุกรายการ \\
เอกสารสัญญา API & รหัส 200 เอกสาร OpenAPI 3.1.0 ระบุรุ่น API 1.1.0 \\
ส่วนหัวความปลอดภัยของคำตอบ & ได้ \code{nosniff} และ \code{no-referrer} ครบทั้งคู่ \\
ผลตรวจของ validator ต่อตารางที่อัปโหลด & ค่า \code{scheduleHash} ตรงกับหลักฐานชุดหลัก สถานะ \code{valid_with_warnings} \\
การสร้างโครงการพร้อมสิทธิ์เจ้าของ & รหัส 201 ค่า \code{snapshotHash} ตรงกับค่าที่คำนวณจากโครงการชุดหลัก \\
การบันทึก revision ใหม่ & รหัส 200 ได้หมายเลข revision ต่างจากตัวตั้ง \\
การบันทึกบน revision ที่ล้าสมัย & รหัส 409 \code{REVISION_CONFLICT} พร้อมระบุ revision ปัจจุบัน \\
การตัดข้อมูลของลิงก์แชร์ขอบเขตตารางสอบ & ไม่เปิดเผย projectId ตัดรายชื่อผู้คุม ปัญหานำเข้า และ artifacts ออก \\
การเก็บหลักฐานของลิงก์แชร์ขอบเขตตรวจสอบ & เข้าถึง audit ได้ คง artifacts ครบ 4 จาก 4 รายการ \\
ความครบของลิงก์แชร์โครงการเต็ม & เปิดเผย projectId ข้อมูลตรงกับโครงการที่บันทึกทุกประการ \\
ขอบเขตสิทธิ์ของผู้ถือลิงก์แชร์แบบอ่านอย่างเดียว & การสร้างลิงก์ การอ่าน audit และการโหลดแหล่งข้อมูล ปฏิเสธ 403 ทั้งหมด \\
ความถูกต้องของแหล่งข้อมูลที่ดึงกลับ & ไบต์ตรงกับที่อัปโหลดครบ 45 ไบต์ \\
การเพิกถอนลิงก์แชร์ & หลังเพิกถอน การอ่านได้รหัส 401 ส่วนลิงก์ที่เหลือยังใช้งานได้ \\
ความแยกไม่ได้ของข้อผิดพลาดการยืนยันตัวตน & เนื้อความ 401 เหมือนกันทั้ง id ที่ไม่รู้จักและความลับที่ผิด เช่นเดียวกับ 403 \\
ขอบเขตข้อมูลนำเข้า & แหล่งข้อมูลเกิน 20 รายการปฏิเสธ 413 ชื่อยาว 250 อักขระถูกตัดเหลือ 200 \\
การจำกัดอัตราการเขียน & ตอบ 429 ครั้งแรกเมื่อคำขอที่ 227 จาก 259 คำขอ ภายใน 19.5 วินาที \\\bottomrule
\end{longtable}}

การตรวจอัตราการเขียนต้องอ่านร่วมกับพฤติกรรมที่ Cloudflare ระบุไว้สำหรับตัวจำกัดชนิดนี้ คือทำงานแบบผ่อนปรนและเข้ากันได้ในท้ายที่สุด คำขอที่มาถึงพร้อมกันตรวจจากตัวนับที่เก็บไว้ในเครื่อง ทำให้การส่งคำขอเป็นชุดสั้น ๆ ทะลุเพดานได้ การตรวจจึงกดคำขอต่อเนื่องแทน และพบการปฏิเสธ 429 ครั้งแรกเมื่อคำขอที่ 227 รวม 32 ครั้งก่อนหยุดตามงบที่ตั้งไว้ ผลนี้ยืนยันว่าตัวจำกัดใช้งานจริงบนระบบที่เผยแพร่ แต่ไม่ใช่การควบคุมที่แม่นยำตรงเพดาน 120 คำขอต่อนาที

นอกจากการตรวจ API ยังตรวจการไหลของผู้ใช้ผ่านเบราว์เซอร์บนโดเมนเดียวกัน คือนำเข้าชุดข้อมูลหลัก จัดตารางครบทั้งกลางภาคและปลายภาค บันทึกโครงการขึ้นคลาวด์ สร้างลิงก์แชร์ขอบเขตตารางสอบ และเพิกถอนลิงก์เดิม แล้วเปิดลิงก์ที่สร้างไว้ในบริบทเบราว์เซอร์ใหม่ หน้าจอเข้าสู่มุมมองที่แชร์โดยตรง แสดงตารางเดียวกับของเจ้าของแต่เป็นแบบอ่านอย่างเดียว และไม่เปิดเผยข้อมูลนอกขอบเขต เช่น รายชื่อผู้คุม รายละเอียดการบันทึกภาพอยู่ในแฟ้ม \code{docs/screenshots/cloud-screenshots.json} ซึ่งแยกจากชุดภาพวันที่ 5 กันยายน

\begin{figure}[tbp]\centering
\includegraphics[width=\textwidth]{docs/screenshots/cloud-save-share.png}
\caption[หน้าจอบันทึกและแชร์บนระบบที่เผยแพร่]{หน้าจอบันทึกและแชร์บนระบบที่เผยแพร่จริง แสดงโครงการที่บันทึกขึ้นคลาวด์ สิทธิ์เจ้าของ และลิงก์แชร์หนึ่งรายการที่เพิกถอนแล้วกับอีกหนึ่งรายการที่ยังใช้งานได้}
\label{fig:cloud-save}
\end{figure}
\begin{figure}[tbp]\centering
\includegraphics[width=\textwidth]{docs/screenshots/cloud-share-viewer.png}
\caption[มุมมองที่แชร์บนระบบที่เผยแพร่]{มุมมองที่แชร์เมื่อเปิดลิงก์ในบริบทเบราว์เซอร์ใหม่ แสดงแถบอ่านอย่างเดียวพร้อมหมายเลข revision และผลตามประเภทสอบ โดยซ่อนข้อมูลนอกขอบเขตออกจากหน้าจอ}
\label{fig:cloud-view}
\end{figure}

แถวข้อมูลที่การตรวจสร้างยังอยู่ในที่เก็บของระบบที่เผยแพร่ เพราะ API รุ่น 1.1.0 ไม่มีปลายทางลบโครงการ การตรวจอัตราในรอบสุดท้ายสร้างโครงการสำเร็จ 227 รายการ ทุกแถวตั้งชื่อให้ระบุที่มาได้และเป็นข้อมูลสังเคราะห์ทั้งหมด ส่วนความลับของสิทธิ์เจ้าของและลิงก์แชร์ไม่ถูกบันทึกลงแฟ้มหลักฐานใด ใช้เฉพาะชั่วคราวและอ้างถึงเพียงคำนำหน้าแฮชแปดอักขระ

ขอบเขตที่การตรวจนี้ยังไม่ยืนยัน ได้แก่ พฤติกรรมระยะยาว เช่น การหมดอายุของลิงก์แชร์ การทำงานภายใต้โหลดพร้อมกันสูงจากผู้ใช้หลายราย และแผนสำรองหรือกู้คืนข้อมูล การตรวจจึงยืนยันการทำงานตามสัญญา API ไม่ใช่คำรับประกันความพร้อมใช้ของบริการ
\chapter{ผลการทดสอบเชิงหน้าที่}\label{ch:function}
\subsection{ผลการทำงานตามสถานการณ์}
การยืนยันผลทั้งแปดสถานการณ์ผ่านผลคาดหมายที่กำหนดไว้ล่วงหน้า ชุด clean และ stale rule จัดได้ 82 รายการ ส่วนข้อขัดแย้งสอบพร้อมกันและชุดรวมจัดได้ 78 รายการ อีกสี่รายการไม่มีเวลา

\noindent\textbf{หมายเหตุ:} คำว่า ผ่าน ในตารางนี้หมายถึงพฤติกรรมตรงกับกรณีทดสอบ รวมถึงกรณีที่ควรออกด้วยรหัสออก 2 ไม่ได้หมายความว่าชุดข้อมูลผิดพลาดให้ตารางที่พร้อมใช้

\begin{center}\begin{minipage}{\textwidth}\centering
\captionof{table}{ผลการทดสอบทั้งแปดสถานการณ์}\label{tab:scenarios}
\begin{tabular}{L{4.6cm}ccc}\toprule
สถานการณ์ & รหัสออก & จัดได้ & ระยะสุดท้าย \\\midrule
clean & 0 & 82 & สร้าง audit \\
stale rule & 0 & 82 & สร้าง audit \\
same-time contradiction & 2 & 78 & สร้าง audit \\
combined & 2 & 78 & สร้าง audit \\
รหัสวิชาผิดรูปแบบ & 2 & ไม่ค้นหา & ตรวจนำเข้า \\
ไม่มีกลุ่มนักศึกษา & 2 & ไม่ค้นหา & ตรวจนำเข้า \\
เวลาผิดรูปแบบ & 2 & ไม่ค้นหา & ตรวจนำเข้า \\
ชื่อห้องซ้ำ & 2 & ไม่ค้นหา & ตรวจนำเข้า \\\bottomrule
\end{tabular}\end{minipage}\end{center}
หน้าจอภาพรวมสรุปผลในมุมมองเดียว คือจำนวนที่จัดได้ตามประเภทสอบ จำนวนข้อค้นพบ และสถานะรวม ส่วนขอบเขตความครอบคลุม 78 นับเฉพาะรายการที่ค้นหาใหม่ ไม่รวมสี่รายการที่มีเวลานำเข้าแล้ว ผลของสถานการณ์ข้อขัดแย้งสอดคล้องกับขนาดผลกระทบที่อธิบายไว้ในบทที่ \ref{ch:valid} คือเกิดจากคู่วิชาเดียวที่มีสองประเภทสอบ
\subsection{คุณภาพของตารางจากชุดหลัก}
ตัวชี้วัดคุณภาพของตารางที่ได้จากชุดหลักสรุปได้ดังนี้ ตัวเลขทั้งหมดมาจากเขตข้อมูล quality ของผลลัพธ์ที่บันทึกไว้

\begin{figure}[tbp]\centering
\includegraphics[width=0.78\textwidth]{docs/screenshots/overview.png}
\caption[หน้าจอภาพรวมจากชุดข้อมูลหลัก]{หน้าจอภาพรวมจากชุดข้อมูลหลัก แสดงจำนวนที่จัดได้ 41 กลางภาค และ 41 ปลายภาค ข้อค้นพบ และขอบเขตการตรวจสอบ}
\label{fig:overview}
\end{figure}
\begin{center}\begin{minipage}{\textwidth}\centering\small\renewcommand{\arraystretch}{1.2}
\captionof{table}{คุณภาพตารางสอบจากชุดข้อมูลหลัก}\label{tab:quality}
\begin{tabular}{L{8.6cm}c}\toprule
ตัวชี้วัด & ค่า \\\midrule
รายการสอบ managed ที่มีเวลา & 82 \\
คู่รายการที่กลุ่มเดียวกันสอบวันเดียวกัน (sameDayPairs) & 76 \\
คู่รายการที่กลุ่มเดียวกันสอบวันติดกัน (consecutiveDayPairs) & 248 \\
กลุ่มนักศึกษาที่มีสอบวันเดียวกันอย่างน้อยหนึ่งคู่ (groupsWithSameDayExams) & 4 \\
จำนวนสอบพร้อมกันสูงสุดในช่วงเวลาเดียว (peakConcurrentExams) & 3 \\
รายการในวันเสาร์อาทิตย์ & 3 \\
รายการในวันหยุด & 1 \\
จำนวนวันที่ใช้สอบ กลางภาค และปลายภาค & 12 และ 13 \\\bottomrule
\end{tabular}\end{minipage}\end{center}
จำนวนวันนับจากผลที่มีจริง ไม่ใช่จำนวนวันธรรมดาในช่วงที่ตั้งไว้ ตัวชี้วัดสองตัวแรกคือตัวเลขเดียวกับที่หน้าจอตรวจสอบแสดง การอ่านค่าเหล่านี้ต้องจำไว้ว่าภาระสอดคล้องเชิงคู่ ไม่ใช่การนับนักศึกษา และความครอบคลุมที่จัดได้ครบไม่แปลว่าค่าคุณภาพทุกตัวต่ำที่สุดเท่าที่เป็นไปได้
\subsection{ตัวอย่างช่วงสอบหนึ่งช่วงจากผลจริง}
เพื่อให้เห็นรูปธรรมของผลลัพธ์ ตารางนี้แสดงรายการสอบในช่วงบ่ายของวันแรกของช่วงสอบกลางภาค (17 สิงหาคม 2569 เวลา 13:00 ถึง 16:00) ตามผลที่บันทึกใน \code{results-clean.json}

\begin{center}\begin{minipage}{\textwidth}\centering\small\renewcommand{\arraystretch}{1.2}
\captionof{table}{ตัวอย่างรายการสอบที่ใช้ช่วงเวลาเดียวกัน}\label{tab:example}
\begin{tabular}{L{2.6cm}L{2.0cm}L{4.6cm}}\toprule
รหัสวิชา & ประเภท & กลุ่มนักศึกษา \\\midrule
\code{060321085} & กลางภาค & \code{CNS-2R-DE-RA}, \code{IAM-2R-DE-RA}, \code{ITS-2R-DE-RA}, \code{LOG-2R-DE-RA} \\
\code{060321155} & กลางภาค & \code{ITS-3R-DE-RA} \\
\code{060321260} & กลางภาค & \code{IAM-1R-DE-RB} \\\bottomrule
\end{tabular}\end{minipage}\end{center}
สามรายการนี้แชร์ช่วงเวลาเดียวกันได้ เพราะไม่มีกลุ่มนักศึกษาร่วมกันเลย และเป็นช่วงที่มีจำนวนรายการสูงสุดของตารางนี้
\chapter{ผลด้านเวลาและทรัพยากรการค้นหา}\label{ch:perf}
\subsection{วิธีวัดและสภาพแวดล้อม}
การวัดแบบ CLI เริ่มก่อนสร้างโปรเซสและสิ้นสุดเมื่อโปรเซสออกหลังเขียนผล จึงรวมการเริ่ม Node การอ่านแฟ้ม การนำเข้า การจัดตาราง การตรวจ และการเขียน JSON ไม่มีรอบอุ่นเครื่อง เพราะต้นทุนเริ่มโปรเซสเป็นส่วนของสิ่งที่ต้องการวัด

การวัดภายในโปรเซสนำเข้าแต่ละชุดครั้งเดียว เรียก \code{scheduleProject} หนึ่งครั้งโดยไม่วัด แล้ววัดอีกยี่สิบครั้ง ฟังก์ชันนี้รวมการทำสำเนา การสร้างรายการ การค้นหา คุณภาพ และการตรวจ จึงเรียกผลนี้ว่าเวลา scheduleProject ไม่ใช่เวลาของ solver ล้วน ทั้งสองชุดตัวเลขเกิดจาก \code{docs/data/bench.ts} แต่ละขนาดเก็บยี่สิบตัวอย่างต่อวิธีวัด รายงานค่ามัธยฐานพร้อมควอร์ไทล์ 1 และ 3 เพื่อแสดงการกระจายของค่ากึ่งกลาง 50 เปอร์เซ็นต์ ตัวอย่างดิบทุกค่าเก็บไว้ใน \code{benchmarks.json} การวัดครั้งนี้เกิดขึ้นบนเครื่องที่ยังมีแอปพลิเคชันเดสก์ท็อปอื่นทำงานอยู่เบื้องหลัง ผู้ทดลองไม่ได้เปิดงานอื่นโดยตั้งใจ แต่ระบบบันทึกโหลดเฉลี่ยขณะเริ่มวัดไว้ในตารางด้วย ค่าเวลาทั้งหมดจึงควรอ่านเป็นค่าที่รวมสภาพแวดล้อมของเครื่องในขณะนั้น ไม่ใช่ขอบเขตล่างสุดของระบบ

% Generated by docs/data/bench.ts from benchmarks.json; edit bench.ts, not this file.
\begin{center}\begin{minipage}{\textwidth}\centering\captionof{table}{สภาพแวดล้อมการวัด}\label{tab:env}{\small\begin{tabular}{>{\raggedright\arraybackslash}p{3cm}>{\raggedright\arraybackslash}p{9cm}}\toprule Environment B1 & ค่าที่บันทึก \\\midrule
Apple M4 & RAM 16 GB, macOS 26.5.2, arm64 \\
Node & v26.0.0, TZ Asia/Bangkok \\
พลังงาน & AC Power, แบตเตอรี่ 100\%, low power mode ปิด \\
บันทึกเมื่อ & 2026-09-06T14:34:02.430Z (UTC) \\
โหลดเฉลี่ยของระบบขณะเริ่มวัด & 9.78, 12.18, 10.15 (1, 5, 15 นาที; เครื่อง 10 logical core) \\
งานร่วม & ไม่มีงานอื่นที่ผู้ทดลองตั้งใจให้แข่งขันทรัพยากรระหว่างวัด แต่มีแอปพลิเคชันเดสก์ท็อปทำงานอยู่ จึงบันทึกโหลดระบบประกอบ \\\bottomrule\end{tabular}}\end{minipage}\end{center}
\needspace{12\baselineskip}\begin{center}\begin{minipage}{\textwidth}\centering\captionof{table}{เวลาที่วัดได้แยกตามวิธีวัด (หน่วย ms)}\label{tab:runtime}{\small\tabcolsep=2pt\begin{tabular}{>{\raggedright\arraybackslash}p{2.2cm}>{\raggedright\arraybackslash}p{2.9cm}rrrrr}\toprule ชุด & วิธี & มัธยฐาน & ควอร์ไทล์ 1 & ควอร์ไทล์ 3 & ต่ำสุด & สูงสุด \\\midrule
เล็ก (41 วิชา) & CLI & 557.70 & 460.15 & 638.94 & 303.45 & 679.45 \\
เล็ก (41 วิชา) & {\small scheduleProject} & 84.93 & 64.03 & 94.19 & 58.22 & 129.64 \\
กลาง (100 วิชา) & CLI & 986.41 & 918.85 & 1079.09 & 782.78 & 1184.04 \\
กลาง (100 วิชา) & {\small scheduleProject} & 293.30 & 243.16 & 307.21 & 216.80 & 374.37 \\
ใหญ่ (189 วิชา) & CLI & 2121.66 & 1989.26 & 2228.98 & 1666.86 & 2503.62 \\
ใหญ่ (189 วิชา) & {\small scheduleProject} & 848.19 & 744.51 & 935.01 & 639.76 & 1289.99 \\
\bottomrule\end{tabular}}\end{minipage}\end{center}
\needspace{15\baselineskip}\begin{center}\begin{minipage}{\textwidth}\centering\captionof{table}{สถิติการค้นหาและสถานะการพิสูจน์ของแต่ละขนาดชุดข้อมูล}\label{tab:search}{\small\tabcolsep=3pt\begin{tabular}{>{\raggedright\arraybackslash}p{6cm}ccc}\toprule ตัวชี้วัดของการค้นหา & เล็ก (41 วิชา) & กลาง (100 วิชา) & ใหญ่ (189 วิชา) \\\midrule
visitedNodes & 3255 & 194 & 12000 \\
improvementEvaluations & 14326 & 19806 & 8000 \\
feasibilityBudget & 12000 & 12000 & 12000 \\
งบรวมที่ใช้ (เปอร์เซ็นต์ของ 20000) & 87.91 & 100 & 100 \\
coverageOptimal & จริง & จริง & จริง \\
avoidedDaysOptimal & จริง & จริง & เท็จ \\
budgetExhausted & เท็จ & เท็จ & เท็จ \\
improvementBudgetExhausted & เท็จ & จริง & จริง \\
\bottomrule\end{tabular}}\end{minipage}\end{center}

\subsection{การอ่านผลเวลา}
\begin{figure}[tbp]\centering
\begin{tikzpicture}
\begin{axis}[width=13.2cm,height=6.4cm,ylabel={เวลา (ms)},xlabel={จำนวนวิชา managed},
boxplot/draw direction=y,
xtick={1,2,3},xticklabels={41,100,189},ymode=log]
\addplot+[boxplot] table[y=small]{docs/data/charts/runtime-box-cli.dat};
\addplot+[boxplot] table[y=medium]{docs/data/charts/runtime-box-cli.dat};
\addplot+[boxplot] table[y=large]{docs/data/charts/runtime-box-cli.dat};
\end{axis}\end{tikzpicture}
\caption[การกระจายเวลาผ่าน CLI]{การกระจายเวลาแบบครบกระบวนการผ่าน CLI ยี่สิบครั้งต่อขนาด มาตราส่วนแกนตั้งเป็นลอการิทึม กรอบคือพิสัยควอร์ไทล์ เส้นกลางคือมัธยฐาน และจุดที่พล็อตแยกนอกหนวดคือค่าผิดปกติตามเกณฑ์ 1.5 เท่าของพิสัยควอร์ไทล์ แสดงครบทุกค่าที่วัด}
\label{fig:boxcli}
\end{figure}
\begin{figure}[tbp]\centering
\begin{tikzpicture}
\begin{axis}[width=13.2cm,height=6.4cm,ylabel={เวลา (ms)},xlabel={จำนวนวิชา managed},
boxplot/draw direction=y,
xtick={1,2,3},xticklabels={41,100,189},ymode=log]
\addplot+[boxplot] table[y=small]{docs/data/charts/runtime-box-scheduleProject.dat};
\addplot+[boxplot] table[y=medium]{docs/data/charts/runtime-box-scheduleProject.dat};
\addplot+[boxplot] table[y=large]{docs/data/charts/runtime-box-scheduleProject.dat};
\end{axis}\end{tikzpicture}
\caption[การกระจายเวลาเรียก scheduleProject]{การกระจายเวลาเรียก scheduleProject ภายในโปรเซส ยี่สิบครั้งต่อขนาด มาตราส่วนแกนตั้งเป็นลอการิทึม และหนวดกับจุดนอกหนวดใช้เกณฑ์เดียวกับรูปที่ \ref{fig:boxcli}}
\label{fig:boxsched}
\end{figure}
ตัวเลขเวลาเหล่านี้อธิบายชุดวัดเวลาสังเคราะห์สามชุดที่เป็นไปได้ ไม่ใช่กฎการขยายตัวบนการกระจายแบบไม่เอนเอียงของปัญหาการจัดตาราง กล่องทั้งสามในแต่ละรูปเป็นหมวดหมู่ต่างกัน ไม่ใช่จุดบนเส้นแนวโน้ม อ่านค่ามัธยฐานได้ว่า เวลา CLI เพิ่มจากราว 0.56 วินาที เป็นราว 0.99 และ 2.12 วินาที ตามลำดับขนาด 41 100 และ 189 วิชา และเวลาภายในโปรเซสเพิ่มจากราว 0.085 เป็น 0.293 และ 0.848 วินาที การอ่านนี้ไม่ใช่การอนุมานความสัมพันธ์เชิงเส้น หรือความซับซ้อนของอัลกอริทึมจากจำนวนวิชาสามค่า ระดับการเชื่อมโยงของกราฟข้อขัดแย้ง วันเวลาคงที่ กฎสอบพร้อมกัน และช่องวันที่ใช้ได้ล้วนเปลี่ยนปริมาณงานตามโครงสร้างของแต่ละชุดดังตารางโครงสร้างในบทที่ \ref{ch:method} ความแตกต่างระหว่างตัวอย่างสะท้อนความแปรปรวนของการวัดในสภาพแวดล้อมนี้ การทดลองไม่ได้แยกตัวแปรเพื่อพิสูจน์สาเหตุของความแปรปรวน และผลการวัดนี้ไม่เพียงพอสำหรับสรุปแนวโน้มการเปลี่ยนแปลงด้านประสิทธิภาพระหว่างรุ่นของซอฟต์แวร์
\subsection{งานค้นหาและสถานะการพิสูจน์}
ตารางสถิติการค้นหา (ตารางที่ \ref{tab:search}) แสดงงบที่ใช้จริง ชุดเล็กใช้ประมาณ 88 เปอร์เซ็นต์ของงบรวมและพิสูจน์ได้ทั้งความครอบคลุมและจำนวนรายการในวันที่ควรหลีกเลี่ยง ชุดกลางค้นหาพบคำตอบเร็ว และใช้งบส่วนใหญ่ในการปรับปรุงคุณภาพ ชุดใหญ่ใช้งบความเป็นไปได้ครบ 12000 ที่ได้รับ และงบปรับปรุงครบ 8000 พิสูจน์ความครอบคลุมได้ แต่ไม่พิสูจน์ว่าจำนวนรายการในวันที่ควรหลีกเลี่ยงต่ำที่สุด จึงเป็นตัวอย่างที่ความครอบคลุมพิสูจน์ได้แต่คุณภาพด้านวันยังมีช่องให้ปรับปรุง ในชุดใหญ่ \code{visitedNodes} มีค่าเท่ากับ \code{feasibilityBudget} พอดี แต่ \code{budgetExhausted} เป็นเท็จ เนื่องจากสถานะนี้ใช้ระบุเหตุที่ยุติการค้นหา กล่าวคือการยุติโดยยังไม่มีข้อพิสูจน์ ไม่ใช่การที่ตัวนับมีค่าเท่ากับงบ การค้นหาของทุกองค์ประกอบยุติด้วยข้อพิสูจน์หรือการค้นครบก่อนที่งบจะตัดสาขา จึงไม่มีองค์ประกอบใดถูกบันทึกว่างบหมด

\noindent\code{improvementBudgetExhausted} เป็นจริงเพราะการปรับปรุงคุณภาพทำงานจนครบงบที่เหลือ การใช้งบสูงไม่ได้แปลว่าข้อมูลมีข้อจำกัดเสมอไป เพราะตัวแก้ปัญหาอาจถึงขอบเขตที่พิสูจน์ได้ก่อนใช้งบหมด และในทางกลับกันงบหมดโดยไม่มีข้อพิสูจน์หมายเพียงว่าไม่พบคำตอบภายในงบตามที่อธิบายไว้ในบทที่ \ref{ch:solver}
\subsection{ความต้องการห้องรายช่วงเวลา}
\begin{figure}[tbp]\centering
\begin{tikzpicture}
\begin{axis}[width=12.4cm,height=3.0cm,view={0}{90},enlargelimits=false,axis on top,
xlabel={วันสอบ (เดือน-วัน)},
xtick=data,xticklabels from table={docs/data/charts/room-demand-labels.dat}{label},
x tick label style={rotate=45,anchor=east,font=\scriptsize},
ytick={0,1},yticklabels={เช้า,บ่าย},y dir=reverse,
colormap name=demand,colorbar,colorbar style={width=1.8cm}]
\addplot[matrix plot*,mesh/rows=2,point meta=explicit] table[x=x,y=y,meta=total]{docs/data/charts/room-demand-matrix.dat};
\end{axis}\end{tikzpicture}
\caption[ความต้องการห้องรายวันและช่วงเวลา]{จำนวนรายการสอบ managed ต่อวันและช่วงเวลาของชุดหลัก สีเข้มคือความต้องการห้องพร้อมกันสูง ช่วงว่างระหว่างเดือนสองคือช่วงระหว่างกลางภาคกับปลายภาค}
\label{fig:heatmap}
\end{figure}
รูปนี้ตอบคำถามปฏิบัติว่าความต้องการห้องกระจุกตัวอยู่ที่ใดของปฏิทิน ค่าที่แสดงคือความต้องการ จำนวนสูงสุดต่อช่วงคือ 3 ซึ่งสอดคล้องกับค่า peakConcurrentExams ของตารางคุณภาพในบทที่ \ref{ch:function} รายการเต็มวันถูกนับทั้งช่วงเช้าและบ่าย จึงห้ามรวมทุกช่องแล้วเรียกว่าจำนวนรายการสอบไม่ซ้ำ การแยกความต้องการห้องคอมพิวเตอร์ใช้ชื่อห้องเรียนต้นทางที่ขึ้นต้น C1 ซึ่งปรากฏเก้าช่วงในข้อมูลเดียวกัน
\subsection{ความอ่อนไหวของพารามิเตอร์ของฟังก์ชันต้นทุน}
ค่าน้ำหนักภาระกลุ่ม 10 และ 3 และสัดส่วนงบความเป็นไปได้ 0.6 เป็นพารามิเตอร์ที่ผู้พัฒนาเลือก ไม่ใช่ค่าที่ทฤษฎีบังคับ การทดสอบความอ่อนไหวจึงเปลี่ยนค่าเหล่านี้ในสำเนาชั่วคราวของซอร์สโค้ดโดยไม่แตะต้อง working tree ที่บันทึกไว้ ผลทุกชุดบันทึกใน \code{docs/data/sensitivity.json} พร้อมแฮชของซอร์สที่แก้ไข โดยชุดฐานต้องให้ผลตรงกับค่าที่บันทึกไว้ทุกประการก่อนจึงอ่านค่าอื่น

\begin{center}\begin{minipage}{\textwidth}\centering
\captionof{table}[ผลทดสอบความอ่อนไหวของพารามิเตอร์]{ผลทดสอบความอ่อนไหวต่อน้ำหนักภาระกลุ่มและสัดส่วนงบความเป็นไปได้ จำนวนรายการที่จัดได้เท่ากันทุกการตั้งค่า คือ 82, 200 และ 378 รายการตามลำดับ}\label{tab:sensitivity}
{\small\renewcommand{\arraystretch}{1.15}\begin{tabular}{L{3.6cm}cccc}\toprule
 & \multicolumn{2}{c}{ชุดเล็ก} & \multicolumn{2}{c}{ชุดกลาง} \\\cmidrule(lr){2-3}\cmidrule(lr){4-5}
ค่าที่เปลี่ยน & วันเดียวกัน & วันติดกัน & วันเดียวกัน & วันติดกัน \\\midrule
ค่าเดิม (10 : 3, 0.6) & 76 & 248 & 55 & 265 \\
น้ำหนัก 5 : 1 & 76 & 248 & 54 & 269 \\
น้ำหนัก 3 : 1 & 76 & 248 & 54 & 265 \\
น้ำหนัก 1 : 0 & 76 & 287 & 54 & 351 \\
งบความเป็นไปได้ 0.4 & 76 & 248 & 55 & 265 \\
งบความเป็นไปได้ 0.8 & 76 & 248 & 55 & 265 \\\bottomrule
\end{tabular}}\end{minipage}\end{center}
ผลปรากฏว่าจำนวนรายการที่จัดได้คงเดิมในทุกค่าที่ทดลองภายในสามชุดสังเคราะห์นี้ จึงไม่พบความไวสูงของความครอบคลุมต่อค่าพารามิเตอร์ในช่วงที่ทดลอง รายละเอียดที่ต่างกันอยู่ที่การกระจายวัน กล่าวคือ อัตราส่วน 5 : 1 และ 3 : 1 เปลี่ยนจำนวนคู่วันเดียวกันและวันติดกันจากค่าเดิมไม่เกินราวสองเปอร์เซ็นต์ในสองชุดแรก ขณะที่การไม่สนใจวันติดกันเลย (1 : 0) ทำให้คู่วันติดกันเพิ่มขึ้นราว 16 เปอร์เซ็นต์ในชุดเล็ก และราว 32 เปอร์เซ็นต์ในชุดกลาง โดยรายการสอบกลับมากระจุกตัวในวันเดียวกันแทน การเปลี่ยนสัดส่วนงบระหว่างค้นหากับปรับปรุงเพียงย้ายงบโดยไม่เปลี่ยนผลลัพธ์ของชุดทดลองนี้ ข้อสรุปคือค่าที่เลือกไว้ทำให้ระบบให้น้ำหนักกับการมีสอบหลายรายการในวันเดียวกันมากกว่าการมีสอบในวันติดกันตามเจตนา
\chapter{บทสรุป ข้อจำกัด และงานต่อไป}\label{ch:conclusion}
\subsection{สิ่งที่หลักฐานยืนยัน}
ตารางนี้ตอบคำถามประเมินผล E1 ถึง E5 ตามที่ตั้งไว้ในบทที่ 1 โดยสรุปสถานะที่พิสูจน์ได้ของระบบเมื่อสิ้นสุดรายงาน

\begin{center}\begin{minipage}{\textwidth}\centering
\captionof{table}[สถานะหลักฐานเมื่อสิ้นสุดรายงาน]{สถานะหลักฐานของระบบเมื่อสิ้นสุดรายงาน จัดตามคำถามประเมินผล E1 ถึง E5}\label{tab:conclusion}
\begin{tabular}{L{4.4cm}L{3.2cm}L{5.4cm}}\toprule
ด้าน & สถานะหลักฐาน & รายละเอียด \\\midrule
E1 การจัดวันเวลา (82 จาก 82) & ยืนยันแล้ว & ผ่านการตรวจซ้ำด้วยเส้นทางตรวจที่แยกจากตัวค้นหา \\
E2 ข้อกำหนดหลัก & ยืนยันแล้ว & ไม่พบกลุ่มนักศึกษาทับซ้อน กฎสอบพร้อมกันมีวันเวลาเดียวกัน และรักษาวันเวลาต้นทาง \\
E3 จำนวนจัดได้ถึงขอบเขตบนที่พิสูจน์ได้ & ยืนยันแล้ว & ทั้งสามขนาดของชุดวัดเวลา \\
E4 ความเหมาะสมของวันที่ควรหลีกเลี่ยง & ยืนยันบางส่วน & ยืนยันในชุดเล็กและกลาง \\
E5 ส่วนติดต่อผู้ใช้และ cloud & ยืนยันแล้ว & มุมมองหลักเจ็ดมุมมอง และการบันทึกกับแบ่งปันโครงการบนระบบที่เผยแพร่จริงตามหัวข้อ \ref{sec:cloud} การโต้ตอบที่ลึกกว่าการเปิดดูยังเป็นงานต่อไป \\\bottomrule
\end{tabular}\end{minipage}\end{center}
การทดสอบด้วยแฟ้มชุดข้อมูลสังเคราะห์ยืนยันว่าระบบรักษาวันเวลาที่นำเข้า จัดวิชาตามกฎสอบพร้อมกัน หลีกเลี่ยงการทับซ้อนของกลุ่มที่กำหนด และแสดงผลกรณีผิดพลาดโดยแยกการปฏิเสธข้อมูลกับข้อขัดแย้งการจัดสอบได้ เครื่องมือทดสอบตรวจทั้งจำนวนและพฤติกรรม จึงทำให้ข้อสรุปตรวจซ้ำได้โดยไม่ต้องอาศัยการอ่านภาพตารางเพียงอย่างเดียว การค้นหาให้ความครอบคลุมมาก่อน แล้วจึงลดการใช้วันที่หลีกเลี่ยง ภาระกลุ่ม และความต้องการใช้ทรัพยากรพร้อมกันในช่วงเวลาเดียวกันตามลำดับ วิธีนี้ทำให้ผู้ใช้ตีความผลได้ตามลำดับเป้าหมาย และการทดสอบความอ่อนไหวพบว่าภายในช่วงค่าที่ทดลอง ไม่พบความไวสูงของความครอบคลุมต่อค่าพารามิเตอร์ดังกล่าว
\subsection{ข้อจำกัดที่ต้องคงไว้ในผลสรุป}
ข้อจำกัดหลักที่ผู้อ่านต้องคงไว้เมื่อตีความผลสรุปรวบรวมไว้ในตารางนี้ รายละเอียดของแต่ละข้ออยู่ในบทที่อ้างถึงตลอดรายงาน

\begin{center}\begin{minipage}{\textwidth}\centering
\captionof{table}[ข้อจำกัดของผลการทดสอบ]{ข้อจำกัดของผลการทดสอบและผลต่อการตีความ}\label{tab:limits}
\begin{tabularx}{\textwidth}{@{}L{0.44\textwidth}Y@{}}\toprule
ข้อจำกัด & ผลต่อการตีความ \\\midrule
ชุดข้อมูลเป็นสังเคราะห์ และชุดใหญ่เลือกกรณีที่จัดครบได้ & ไม่ใช้แทนการกระจายความยากของข้อมูลจริงทั้งหมด และไม่ครอบคลุมกรณีที่จัดไม่ได้ \\
ความเหมาะสมของวันที่ควรหลีกเลี่ยงยังไม่พิสูจน์ในชุดใหญ่ & ชุดใหญ่มีข้อพิสูจน์ความครอบคลุม แต่ยังไม่ใช่ข้อพิสูจน์คุณภาพทุกด้าน \\
การวัดเวลารันบนเครื่องเดียวภายใต้โหลดที่บันทึกไว้ & ห้ามตีความเป็นคำรับประกันสมรรถนะ \\
การโต้ตอบกับส่วนติดต่อผู้ใช้ทดสอบเฉพาะมุมมองหลัก & การแก้ตารางด้วยมือและขั้นตอนที่ลึกกว่ายังไม่ได้รับการยืนยัน \\
การตรวจระบบที่เผยแพร่เป็นการใช้งานระดับต่ำในวันเดียว & ไม่ครอบคลุมพฤติกรรมระยะยาว โหลดพร้อมกันสูง และแผนสำรองหรือกู้คืนข้อมูล \\\bottomrule
\end{tabularx}\end{minipage}\end{center}
\subsection{งานต่อไป}
งานถัดไปด้านรายงานคือทดสอบการโต้ตอบที่ลึกขึ้นในเบราว์เซอร์ เช่น การแก้ตารางด้วยมือ และการจัดการข้อขัดแย้งของ revision ที่หน้าจอ โดยภาพแต่ละภาพควรระบุภาษา ขนาดหน้าจอ และแฮชข้อมูลควบคู่กันตามแบบแฟ้มหลักฐานปัจจุบัน

งานถัดไปด้านระบบคือประเมินการค้นหาด้วยชุดที่เพิ่มความหนาแน่นของกราฟอย่างควบคุม แทนการเพิ่มจำนวนวิชาอย่างเดียว เพื่อแยกผลของขนาดกับโครงสร้างข้อกำหนด
\subsection{ขั้นตอนทำซ้ำที่แนะนำ}
เริ่มจากตรวจอัตลักษณ์โปรแกรมและสภาพแวดล้อมตาม snapshot จากนั้นสร้าง fixture และเรียกตัวตรวจผล เมื่อทุกเงื่อนไขผ่านจึงหยุดงานที่แข่งขันทรัพยากรและรันวัดเวลา หลังได้หลักฐานครบจึงบันทึกแฮชและสร้างเอกสาร หลักฐานทั้งหมดควรมาจาก revision เดียวกันตลอดกระบวนการ หากซอร์สโค้ดเปลี่ยนหลังสร้าง snapshot ต้องบันทึก snapshot และหลักฐานชุดใหม่แทนการอ้าง snapshot เดิม

\begin{center}\begin{minipage}{\textwidth}\centering
\captionof{table}[ขั้นตอนทำซ้ำที่แนะนำ]{ขั้นตอนทำซ้ำที่แนะนำ: แต่ละขั้นใช้เครื่องมือหลักที่ระบุ}\label{tab:repro}
\begin{tabular}{L{3.6cm}L{9.4cm}}\toprule
ขั้น & เครื่องมือหลัก \\\midrule
สร้างข้อมูล & \cmdcell{node docs/data/generate.ts} \\
สถิติโครงสร้างชุดข้อมูล & \cmdcell{node docs/data/fixture-stats.ts} \\
ยืนยันผล & \cmdcell{node docs/data/verify-results.ts} \\
วัดเวลา & \cmdcell{node docs/data/bench.ts} พร้อมตัวเลือก \cmd{measure} \\
ทดสอบความอ่อนไหว & \cmdcell{node docs/data/sensitivity.ts} \\
ตรวจระบบ cloud ที่เผยแพร่ & \cmdcell{node docs/data/cloud-verification.ts} \\
บันทึกหลักฐาน & \cmdcell{node docs/repro/snapshot.ts evidence} \\
สร้าง PDF & \cmdcell{./docs/build.sh} \\
ตรวจสถานะโปรแกรม & \cmdcell{node docs/repro/snapshot.ts post} \\\bottomrule
\end{tabular}\end{minipage}\end{center}
แฟ้ม \code{checksums.sha256} เก็บแฮช PDF และหลักฐานฉบับสุดท้ายแยกจาก snapshot เพื่อไม่ให้เกิดวงจรที่ PDF ต้องบรรจุแฮชของตัวเอง การเปลี่ยนหลักฐานหลังสร้าง PDF ทำให้ต้องเริ่มใหม่จากขั้นบันทึกหลักฐาน การเก็บซอร์สเอกสารและเครื่องมือไว้พร้อมผลลัพธ์ช่วยให้แก้ไขรอบถัดไปได้โดยไม่ต้องสร้างวิธีทำงานใหม่
\clearpage\phantomsection\addcontentsline{toc}{section}{เอกสารอ้างอิง}\section*{เอกสารอ้างอิง}
\begingroup
\setlength{\parindent}{0pt}
\newcommand{\refentry}{\noindent\hangindent=1.5em\hangafter=1}
\subsection*{เอกสารวิชาการ}
\refentry Russell, S., \& Norvig, P. (2021). \emph{Artificial intelligence: A modern approach} (4th US ed.). ส่วนบท Constraint Satisfaction Problems. \code{https://aima.cs.berkeley.edu/} (เข้าถึง 5 กันยายน 2569)

\refentry Carnegie Mellon University. (n.d.). \emph{15-281: Constraint satisfaction problems} [โน้ตประกอบการสอน]. \code{https://www.cs.cmu.edu/~15281/coursenotes/constraints/} (เข้าถึง 5 กันยายน 2569)
\subsection*{เอกสารและหลักฐานของโครงการ}
\refentry ระบบจัดตารางสอบ [ซอฟต์แวร์]. (2569). ฉบับตาม Git revision ที่บทที่ \ref{ch:identity} ระบุ ภายใต้สัญญาอนุญาต MIT. GitHub. \code{https://github.com/ItsBreadTime/examscheduler}
\endgroup
\appendix
\renewcommand{\thechapter}{\Alph{chapter}}
\setcounter{chapter}{0}
\chapter{สัญญาข้อมูลนำเข้า}\label{app:input}
\subsection{รายวิชาแบบตารางแบน}
รายการแฟ้มที่ภาคผนวกกำหนดขอบเขตมีตารางรายวิชา กฎ ห้อง และ \code{manifest/settings} เท่านั้น แฟ้มชนิดอื่นอยู่นอกขอบเขต คอลัมน์บังคับของรูปแบบตารางแบนคือ \code{courseCode}, \code{courseName}, \code{sectionNumber} และ \code{studentGroups} รหัสวิชาต้องเป็นข้อความเก้าหลักเพื่อรักษาศูนย์นำหน้า กลุ่มแยกด้วยขึ้นบรรทัดใหม่ อัฒภาค หรือจุลภาค และผ่านการปรับช่องว่างกับ Unicode NFC

\begin{longtable}{L{5.6cm}L{7.6cm}}\toprule
คอลัมน์เสริม & ความหมาย \\\midrule
\code{midtermDate}, \code{midtermTime} & วันและช่วงสอบกลางภาค \\
\code{finalDate}, \code{finalTime} & วันและช่วงสอบปลายภาค \\
\code{plannedEnrollment} & จำนวนเปิดรับ \\
\code{registeredEnrollment} & จำนวนลงทะเบียน \\
\code{teachingDay}, \code{teachingTime} & วันและเวลาเรียน \\
\code{teachingRoom}, \code{instructors} & ห้องเรียนและข้อมูลผู้สอน \\\bottomrule
\end{longtable}
ชุดรายงานใช้ workbook แบบหัวตารางภาษาไทยที่ตัวนำเข้ารองรับ คอลัมน์วันที่และเวลาเป็นคู่ที่อยู่ติดกัน รวมถึงจำนวนเปิดรับและลงทะเบียน จึงไม่ควรย้ายเฉพาะคอลัมน์ย่อยโดยคงหัวเดิม แถวที่ไม่มีรหัสวิชาสืบต่อวิชาก่อนหน้าในตารางปัจจุบัน หากวันและเวลาสอบว่างทั้งคู่หมายถึงยังไม่ระบุ แต่มีเพียงค่าเดียวถือว่าผิดรูปแบบ
\subsection{กฎและห้อง}
หัวกฎ ไม่มีสอบกลางภาค แปลว่าวิชาในคอลัมน์นั้นถูกยกเว้นจากการจัดสอบกลางภาค ระบบจึงไม่สร้างรายการสอบกลางภาคให้ ส่วนหัว ไม่มีสอบปลายภาค แปลเป็นไม่มีสอบปลายภาค และหัว ไม่มีสอบ ใช้กับทั้งสองประเภท หัวที่มีคำว่าสองช่วงใช้ full day ส่วนหัวที่มีคำว่าตรงกันอ่านรหัสที่เหลือในแถวนั้นเป็นกฎสอบพร้อมกัน

ห้องใช้คอลัมน์แรกชื่อห้อง ตามด้วยแถวสำหรับโซนที่คั่นด้วยช่องว่าง และคอมซึ่งค่าใช่เพิ่มป้าย computer ชื่อห้องซ้ำถูกปฏิเสธ ส่วนความจุที่นั่งไม่ได้กำหนดจากแฟ้มห้องนี้ ตารางเขตข้อมูล manifest อยู่ในหัวข้อถัดไป
\subsection{Manifest และ settings}
\begin{longtable}{L{4.6cm}L{8.6cm}}\toprule
ชื่อเขตข้อมูล & ข้อกำหนด \\\midrule
\code{courseSources} & รายการ path, role และ expectedPrefix เมื่อใช้ตรวจคำนำหน้า \\
\code{rules} & เส้นทางแฟ้มกฎ \\
\code{references} & รายการ path และ kind สำหรับแหล่งห้อง \\
\code{settings.periods} & start/end แยก midterm และ final \\
\code{settings.sessions} & ช่วงนาทีที่ถูกต้องอย่างน้อยหนึ่งช่วง \\
\code{settings.fullDay} & ช่วงนาทีเต็มวัน \\
\code{settings.weekendPolicy} & \code{never}, \code{only_if_necessary}, \code{normal} \\
\code{settings.holidayPolicy} & ค่าชุดเดียวกับ weekendPolicy \\
\code{settings.holidays} & วันที่ ISO ที่ถูกต้อง \\
\code{settings.searchBudget} & จำนวนเต็ม 1 ถึง 1000000 \\
\code{locks} & อยู่ระดับเดียวกับ settings ไม่ใช่สมาชิก SchedulerSettings \\\bottomrule
\end{longtable}
แฟ้มหลักตั้งงบ 20000 ทุกขนาด ช่วงสอบกลางภาคเริ่ม 17 ถึง 28 สิงหาคม 2026 และปลายภาค 14 ถึง 30 ตุลาคม 2026 ช่วงเช้า 540 ถึง 720 นาที ช่วงบ่าย 780 ถึง 960 นาที และเต็มวัน 540 ถึง 960 นาที วันหยุดคือ 20 สิงหาคม 2026 นโยบายทั้งสองเป็น \code{only_if_necessary}
\subsection{การอ่านแฟ้ม}
ตัวอ่านรองรับ xlsx, xls, xlsm, xlsb, csv, tsv, ods และ fods ตามรายการนามสกุลปัจจุบัน CSV และ TSV ต้องเป็น UTF-8 ที่ถูกต้อง แฟ้มว่าง แฟ้มเกิน 25 MiB หรือชนิดไม่รองรับเข้าสู่เส้นทางข้อผิดพลาดการอ่านแฟ้ม ไม่ใช้จำนวน issue เป็นหลักฐานแทนข้อผิดพลาดที่ถูกโยนออกมา

\subsection{การอ่านวันที่และเวลา}\label{app:dates}
\code{parseDate} รับ ISO ที่เป็นวันที่จริง หรือวันที่คั่นด้วยเครื่องหมายทับ ปีสองหลักตีความเป็นปีพุทธศักราช เช่น 69 เพิ่มเป็น 2569 แล้วแปลงเป็น 2026 ปีตั้งแต่ 2400 ขึ้นไปลบ 543 จึงรองรับปีพุทธศักราชเต็มด้วย ไม่ควรใช้ปีสองหลักแล้วคาดหวังความหมายตามคริสต์ศักราชทั่วไป

แฟ้มตารางคำนวณที่เก็บวันที่เป็นตัวเลขและมีรูปแบบวันที่ถูกแปลงตามระบบวันที่ของ workbook ทั้งระบบ 1900 และ 1904 ข้อมูลที่ส่งออกใช้วันแบบ ISO เพื่อให้ไม่กำกวม ส่วนการแสดงผลภาษาไทยสามารถแสดงปีพุทธศักราชโดยไม่เปลี่ยนค่าปฏิทินภายใน

\begin{center}\begin{tabular}{L{3.2cm}L{4.4cm}}\toprule
ข้อมูลนำเข้า & ความหมายที่ใช้ภายใน \\\midrule
\code{17/08/69} & \code{2026-08-17} \\
\code{17/08/2569} & \code{2026-08-17} \\
\code{2026-08-17} & วัน ISO เดิม \\
\code{09:00-12:00} & เริ่ม 540 สิ้นสุด 720 นาที \\\bottomrule
\end{tabular}\end{center}

\chapter{สัญญาข้อมูลผลลัพธ์}\label{app:output}
แฟ้มผลลัพธ์ของ CLI มีโครงสร้างระดับบนเจ็ดส่วน ตารางนี้สรุปขอบเขตของแต่ละส่วนเป็นสัญญาที่ผู้ใช้แฟ้มควรพึ่งพาได้ ขอบเขตของรายงานคือแฟ้มที่ปรากฏใน \code{docs/data/}

\begin{longtable}{L{3.4cm}L{9.8cm}}\toprule
เขตข้อมูล & ความหมาย \\\midrule
\code{project} & สถานะโครงการหลังจัดตาราง รวม events, sections, rooms และ settings \\
\code{unscheduled} & หน่วยที่ไม่ได้รับเวลา พร้อมรหัสรายการ เหตุผล และจำนวนช่องที่ถูกบล็อก \\
\code{search} & สถิติและสถานะการพิสูจน์ เช่น budget, budgetExhausted, coverageUpperBound, coverageOptimal, avoidedDaysOptimal, visitedNodes, improvementEvaluations, maxPlaceableKnown \\
\code{quality} & ตัวชี้วัดคุณภาพ เช่น scheduledManagedEvents, sameDayPairs, consecutiveDayPairs, weekendExams, holidayExams, daysUsed \\
\code{validation} & สถานะรวม overallStatus กับข้อค้นพบและการครอบคลุมของการตรวจ \\
\code{inspection} & ผลตรวจขั้นนำเข้าก่อนการจัดตาราง \\
\code{audit} & หลักฐานสำหรับการตรวจซ้ำ เช่น sourceHash, rulesHash, scheduleHash, เวอร์ชันของ parser และ scheduler และสรุปคุณภาพ \\\bottomrule
\end{longtable}
\enlargethispage{\baselineskip}
สถานะรวมมีค่า \code{valid}, \code{valid_with_warnings} และ \code{invalid} ความหมายของแต่ละค่าอธิบายไว้ในบทที่ \ref{ch:valid} ผู้ใช้แฟ้มควรอ่านสถานะควบคู่กับสถิติการค้นหาเสมอ เพราะการมีแฟ้มผลลัพธ์ไม่ได้แปลว่าตารางสมบูรณ์
\needspace{10\baselineskip}\chapter{แฟ้มหลักฐาน บัญชีรหัส และภาพหน้าจอเพิ่มเติม}\label{app:evidence}
\subsection{รหัสที่สังเกตในชุดทดสอบ}
\begin{longtable}{L{6.6cm}L{6.6cm}}\toprule
รหัส & หมวดตามตำแหน่งข้อมูล \\\midrule
\code{EXAM_ON_WEEKEND}, \code{EXAM_ON_HOLIDAY} & issue การใช้วันตามปฏิทิน \\
\code{STALE_RULE_REFERENCE}, \code{DUPLICATE_RULE} & issue จากกฎ \\
\code{SAME_TIME_CONTRADICTION} & issue ตัวตรวจ และเหตุผลหน่วยที่จัดไม่ได้ แยก namespace \\
\code{UNSCHEDULED_REQUIRED_EXAM} & issue รายการจำเป็นไม่มีเวลา \\
\code{MALFORMED_COURSE_ROW} & issue นำเข้ารายวิชา \\
\code{MISSING_STUDENT_GROUPS} & issue กลุ่มนักศึกษาหาย \\
\code{INVALID_EXAM_TIMING} & issue วันเวลาต้นทาง \\
\code{SOURCE_DATA_INVALID} & issue สถานะข้อมูลที่สร้างรายการ \\
\code{DUPLICATE_ROOM} & issue ห้องซ้ำ \\\bottomrule
\end{longtable}
บัญชีนี้รวมรหัสที่รายงานใช้ในขอบเขตเนื้อหา ไม่ใช่รายการรหัสทั้งหมดที่โปรแกรมอาจผลิต รหัสออกจากโปรเซสเป็นจำนวนเต็ม และแยกจาก issue ส่วนชื่ออย่าง \code{coverageUpperBound} หรือ \code{provablyImpossibleEvents} เป็นเขตสถิติ ไม่ใช่รหัสปัญหา
\subsection{การจัดเก็บหลักฐาน}
\code{run-records.json} เก็บคำสั่ง โฟลเดอร์ทำงาน stdout stderr รหัสออก ระยะที่ไปถึง และตำแหน่งผล \code{verification.json} เก็บผลตรวจและชื่อเงื่อนไขหลัก \code{benchmarks.json} เก็บตัวอย่างดิบ สภาพแวดล้อม และวิธีวัด \code{fixture-stats.json} เก็บสถิติโครงสร้างของแต่ละชุดข้อมูล \code{sensitivity.json} เก็บผลทดสอบความอ่อนไหวพร้อมแฮชของซอร์สที่แก้ไขในสำเนาชั่วคราว \code{cloud-verification.json} เก็บผลตรวจ 16 รายการของสัญญา API บนระบบที่เผยแพร่จริง และ \code{cloud-screenshots.json} เก็บรายละเอียดของภาพหน้าจอชุด cloud ซึ่งบันทึกแยกจากชุดวันที่ 5 กันยายน ส่วน charts/ เป็นตารางที่กราฟอ่านโดยตรง เครื่องมือทดสอบเรียก \cmd{node src/cli.ts} โดยตรง เพื่อไม่ให้ข้อความประกอบของ npm ปะปนกับ JSON ที่บันทึกไว้

คำแปลศัพท์อยู่ในบทที่ 1 และใช้ร่วมกับชื่ออังกฤษเดียวกันตลอดรายงาน ภาพหน้าจอในภาคผนวกนี้ใช้ชุดข้อมูลหลักสำหรับมุมมองกลุ่มนักศึกษาและการตรวจสอบ ส่วนหน้าปัญหาในบทที่ \ref{ch:method} ใช้ชุดกฎอ้างอิงล้าสมัยเพื่อแสดงที่มาของคำเตือน

หน้าจอกลุ่มนักศึกษาให้เลือกกลุ่มเพื่อดูเฉพาะรายการสอบของกลุ่มนั้น เมื่อเลือกกลุ่มแล้ว ตารางจะแสดงรายวิชา วันเวลา และสถานะของรายการ
\begin{figure}[p]\centering
\includegraphics[width=\textwidth]{docs/screenshots/groups.png}
\caption[หน้าจอกลุ่มนักศึกษา]{หน้าจอกลุ่มนักศึกษาจากชุดข้อมูลหลัก แสดงตารางสอบเฉพาะของกลุ่มที่เลือก}
\label{fig:groups}
\end{figure}
หน้าจอตรวจสอบสรุปแฮชของตารางสอบ ข้อมูลต้นทาง และเงื่อนไข พร้อมรายการตรวจสอบและสถานะของแต่ละกลุ่มตรวจ
\begin{figure}[p]\centering
\includegraphics[width=\textwidth]{docs/screenshots/audit.png}
\caption[หน้าจอตรวจสอบ]{หน้าจอตรวจสอบจากชุดข้อมูลหลัก แสดงแฮช รายการตรวจสอบ และคุณภาพตารางสอบ}
\label{fig:audit}
\end{figure}
\subsection{ปฏิทินสอบและนโยบายวัน}
ชุดหลักใช้วันหยุด 20 สิงหาคม 2569 ซึ่งอยู่ในช่วงสอบกลางภาค และมีวันเสาร์อาทิตย์ 22 และ 23 สิงหาคมอยู่ในช่วงเดียวกัน การทดสอบจึงพิสูจน์พฤติกรรมกับวันที่มีอยู่จริงในโดเมน ไม่ใช่ใส่วันหยุดไว้นอกช่วงแล้วถือว่าครอบคลุมกรณีวันหยุด
\begin{figure}[tbp]\centering
\includegraphics[width=\textwidth]{docs/screenshots/calendar.png}
\caption[หน้าจอปฏิทินสอบรายสัปดาห์]{หน้าจอปฏิทินจากชุดข้อมูลหลัก แสดงตารางสอบรายสัปดาห์พร้อมวันที่ถูกกำหนดตายตัว}
\label{fig:calendar}
\end{figure}

หน้าจอปฏิทินแสดงตารางสอบรายสัปดาห์ โดยเริ่มต้นวันอาทิตย์ แต่ละคอลัมน์คือวันในสัปดาห์และแต่ละแถวคือสัปดาห์ที่ต่อเนื่องกัน วันที่ถูกกำหนดตายตัวจะถูกทำเครื่องหมาย และรายการสอบเต็มวันกินทั้งช่วงเช้าและบ่าย

\needspace{24\baselineskip}\chapter{การเชื่อมโยงเนื้อหากับซอร์สโค้ด}\label{app:trace}
หัวข้อในรายงานที่อ้างอิงแฟ้ม implementation รวบรวมไว้ในตารางนี้ รายการที่ระบุคือแฟ้มที่แต่ละหัวข้ออ้างข้อเท็จจริง ไม่ใช่รายการแฟ้มทั้งหมดของโครงการ

\begin{longtable}{L{4.6cm}L{8.6cm}}\toprule
หัวข้อในรายงาน & แฟ้ม implementation \\\midrule
\endfirsthead
\toprule
หัวข้อในรายงาน & แฟ้ม implementation \\\midrule
\endhead
บทที่ 2 องค์ประกอบของระบบและช่องทางเรียกใช้ & \codeS{src/ui/jobs.ts;src/validate.ts;src/lib/api.ts;worker/index.ts} \\
บทที่ 2 การนำเข้าและการสร้างรายการสอบ & \codeS{src/lib/import/spreadsheet.ts;src/lib/project.ts;src/lib/model.ts;src/lib/rules.ts} \\
บทที่ \ref{ch:arch} การแทนวันเวลา และภาคผนวก \ref{app:dates} & \code{src/lib/time.ts}\newline ฟังก์ชัน \code{overlaps} และ \code{timingsOverlap} \\
บทที่ \ref{ch:solver} การค้นหาและการปรับปรุง & \codeS{src/lib/scheduler.ts;src/lib/solver.ts} ฟังก์ชัน \codeS{scheduleProject;solveUnits} \\
บทที่ \ref{ch:valid} การตรวจซ้ำและ audit & \codeS{src/lib/validator.ts;src/lib/audit.ts} ฟังก์ชัน \codeS{validateProject;buildAuditBundle} \\
บทที่ \ref{ch:method} เครื่องมือทดสอบและความคาดหวัง & \codeS{docs/data/verify-results.ts;docs/data/expected.json;docs/data/verification.json} \\
บทที่ \ref{ch:method} การตรวจสอบชั้น cloud บนระบบที่เผยแพร่จริง & \codeS{worker/index.ts;docs/data/cloud-verification.ts} \\
ภาคผนวก A สัญญาข้อมูลนำเข้า & \codeS{src/lib/import/courses.ts;src/lib/rules.ts;src/lib/import/rooms.ts;src/lib/types.ts} \\\bottomrule
\end{longtable}
\end{document}
